% ==============================================================================
% ENTREGABLE 2, BLOQUE D: EXTENSION MULTIMODAL DEL MOTOR DE CONSULTAS
%
% Fragmento para insertar en informe.tex despues de las secciones de los bloques
% A, B y C. Se apoya en el preambulo que ya trae el documento: booktabs,
% tabularx, listings, float, pgfplots con el estilo "figura", el estilo tikz
% "nota" y los colores cHeap, cSeq, cBtree, cHash, cGrid, cEje y cNota.
%
% HACE FALTA UNA LINEA MAS EN EL PREAMBULO:
%
%     \usepackage[T1]{fontenc}
%
% justo despues de \usepackage[utf8]{inputenc}. Sin ella la codificacion de
% fuente por defecto imprime < como invertida de admiracion y > como invertida
% de interrogacion dentro de \texttt y de los listados, de modo que 'id < 10' y
% el operador <=> salen mal. Afecta tambien a lo que ya esta escrito en la
% seccion del planificador del entregable 1 ('id >= 100', 'id <= 500').
%
%   % ==============================================================================
% ENTREGABLE 2, BLOQUE D: EXTENSION MULTIMODAL DEL MOTOR DE CONSULTAS
%
% Fragmento para insertar en informe.tex despues de las secciones de los bloques
% A, B y C. Se apoya en el preambulo que ya trae el documento: booktabs,
% tabularx, listings, float, pgfplots con el estilo "figura", el estilo tikz
% "nota" y los colores cHeap, cSeq, cBtree, cHash, cGrid, cEje y cNota.
%
% HACE FALTA UNA LINEA MAS EN EL PREAMBULO:
%
%     \usepackage[T1]{fontenc}
%
% justo despues de \usepackage[utf8]{inputenc}. Sin ella la codificacion de
% fuente por defecto imprime < como invertida de admiracion y > como invertida
% de interrogacion dentro de \texttt y de los listados, de modo que 'id < 10' y
% el operador <=> salen mal. Afecta tambien a lo que ya esta escrito en la
% seccion del planificador del entregable 1 ('id >= 100', 'id <= 500').
%
%   % ==============================================================================
% ENTREGABLE 2, BLOQUE D: EXTENSION MULTIMODAL DEL MOTOR DE CONSULTAS
%
% Fragmento para insertar en informe.tex despues de las secciones de los bloques
% A, B y C. Se apoya en el preambulo que ya trae el documento: booktabs,
% tabularx, listings, float, pgfplots con el estilo "figura", el estilo tikz
% "nota" y los colores cHeap, cSeq, cBtree, cHash, cGrid, cEje y cNota.
%
% HACE FALTA UNA LINEA MAS EN EL PREAMBULO:
%
%     \usepackage[T1]{fontenc}
%
% justo despues de \usepackage[utf8]{inputenc}. Sin ella la codificacion de
% fuente por defecto imprime < como invertida de admiracion y > como invertida
% de interrogacion dentro de \texttt y de los listados, de modo que 'id < 10' y
% el operador <=> salen mal. Afecta tambien a lo que ya esta escrito en la
% seccion del planificador del entregable 1 ('id >= 100', 'id <= 500').
%
%   % ==============================================================================
% ENTREGABLE 2, BLOQUE D: EXTENSION MULTIMODAL DEL MOTOR DE CONSULTAS
%
% Fragmento para insertar en informe.tex despues de las secciones de los bloques
% A, B y C. Se apoya en el preambulo que ya trae el documento: booktabs,
% tabularx, listings, float, pgfplots con el estilo "figura", el estilo tikz
% "nota" y los colores cHeap, cSeq, cBtree, cHash, cGrid, cEje y cNota.
%
% HACE FALTA UNA LINEA MAS EN EL PREAMBULO:
%
%     \usepackage[T1]{fontenc}
%
% justo despues de \usepackage[utf8]{inputenc}. Sin ella la codificacion de
% fuente por defecto imprime < como invertida de admiracion y > como invertida
% de interrogacion dentro de \texttt y de los listados, de modo que 'id < 10' y
% el operador <=> salen mal. Afecta tambien a lo que ya esta escrito en la
% seccion del planificador del entregable 1 ('id >= 100', 'id <= 500').
%
%   \input{entregable2_motor}
%
% Si los experimentos se agrupan todos en la seccion de experimentacion, la
% subseccion "Experimento 9" se mueve alla sin tocar el resto.
% ==============================================================================

\section{Extensión Multimodal del Motor de Consultas}

En el primer entregable el motor resolvía un solo tipo de pregunta: comparar una
columna contra una constante. Este entregable le suma tres formas de buscar que
no se parecen entre sí ni a la anterior. La espacial pide los puntos cercanos a
una coordenada, la textual ordena documentos por relevancia frente a una
consulta de palabras, y la vectorial ordena imágenes por parecido frente a otra
imagen. Las tres devuelven filas ordenadas por una magnitud que la tabla no
guarda en ninguna columna.

El trabajo de esta parte fue que el motor aceptara las tres sin convertirse en
cuatro motores pegados. Lo que hicimos fue extender la gramática, agregar un
contrato común que cualquier estructura de índice puede implementar, y meter en
el planificador la decisión que el enunciado pide reportar: cuándo conviene
acotar las filas antes de consultar el índice y cuándo conviene lo contrario.
Sobre eso se construyó la aplicación de identificación facial, que es la que
pone los tres niveles a trabajar en una sola consulta.

\subsection{Gramática de las Consultas Multimodales}

\subsubsection{Las tres familias de predicado}

Las extensiones se añadieron al analizador léxico y al parser de descenso
recursivo que ya existían, sin cambiar la precedencia de operadores ni la forma
del árbol sintáctico. Cada familia entra por un sitio distinto de la gramática,
y esa diferencia no es estética: refleja dónde encaja cada una en la semántica
de un \texttt{SELECT}.

\begin{table}[H]
\centering
\caption{Dónde entra cada familia en la gramática}
\begin{tabularx}{\textwidth}{@{}llX@{}}
\toprule
\textbf{Familia} & \textbf{Lugar} & \textbf{Sintaxis} \\
\midrule
Espacial, radio & predicado del \texttt{WHERE} & \texttt{col WITHIN\_RADIUS((lon, lat), r)} \\
Espacial, k-NN & cláusula \texttt{ORDER BY} & \texttt{ORDER BY DISTANCE(col, (lon, lat))} \\
Textual & predicado del \texttt{WHERE} & \texttt{col MATCH 'términos'} \\
Textual, puntaje & lista de proyección & \texttt{COSINE\_SCORE() AS rel} \\
Vectorial & lista de proyección & \texttt{col \textless=\textgreater{} '[v1, ..., vn]' AS sim} \\
\bottomrule
\end{tabularx}
\end{table}

El radio es un predicado porque recorta el conjunto de filas y se puede combinar
con \texttt{AND} como cualquier otra condición. El k-NN, en cambio, no recorta
nada: ordena, y el recorte lo hace el \texttt{LIMIT}. Por eso vive en el
\texttt{ORDER BY} y no en el \texttt{WHERE}, que fue la primera cosa que tuvimos
que decidir y la que más ordenó el resto del diseño.

La distancia vectorial es el caso más incómodo de los cinco. No es un predicado
ni una cláusula de orden: es una columna calculada que no está en la tabla. El
operador \texttt{\textless=\textgreater{}} aparece en la proyección, produce una similitud por fila y
obliga a declarar un alias, porque sin nombre el \texttt{ORDER BY} no tiene cómo
referirse a ella. El motor rechaza la consulta si falta:

\begin{lstlisting}[language=SQL, caption={Declaración y uso del índice vectorial}]
CREATE TABLE rostros (
  id        INT PRIMARY KEY,
  lote      CHAR(8),
  embedding VARCHAR(23)
) USING HEAP;

CREATE INDEX ix_rostros_ivf ON rostros (embedding)
  USING IVF WITH (mode = 'face', nprobe = 4);

SELECT id, lote, embedding <=> '[0.111721, 0.152196, ...]' AS similitud
FROM rostros
WHERE id < 1000
ORDER BY similitud DESC
LIMIT 5;
\end{lstlisting}

\subsubsection{La columna indexada guarda una referencia, no el vector}

La tabla no almacena los 128 flotantes de cada rostro. Guarda la ruta de la
imagen, y es el índice el que resuelve de ahí el vector. Esto no fue una
concesión: es la misma idea con la que el R-Tree del bloque espacial indexa un
\texttt{POINT(lon lat)} escrito en una columna de texto. La tabla guarda la
referencia y la estructura guarda la geometría.

La consecuencia práctica es que un registro de \texttt{rostros} ocupa 36 bytes y
entran 101 por página, así que las 13\,000 filas caben en 129 páginas. Si el
vector viviera en la tabla, cada fila pesaría 512 bytes solo de flotantes y la
misma tabla ocuparía más de 1\,700 páginas, con el agravante de que ninguna
consulta relacional necesita leer esos bytes.

\subsection{Acoplamiento de las Tres Familias de Índice}

\subsubsection{Un contrato que no sabe de geometría ni de texto}

Las tres estructuras nuevas se enchufan por la misma interfaz que ya usaban el
árbol B$^+$ y el hash extensible. Un índice declara qué puede servir y el motor
le pregunta antes de registrarlo; si no puede, devuelve el motivo en texto y ese
motivo es lo que llega al usuario:

\begin{lstlisting}[caption={Motivo de rechazo de un índice}]
CREATE INDEX ix ON rostros (id) USING IVF;

  rostros.id debe guardar la ruta de la imagen en CHAR/VARCHAR
\end{lstlisting}

El despachador mantiene la lista de candidatos y encamina cada operación al que
aceptó ese índice. Lo que no acepta nadie cae a un sustituto en memoria, de modo
que una consulta sigue respondiendo aunque la estructura que le correspondía no
esté disponible. El cliente muestra en qué estructura quedó cada índice, y en
las mediciones de este informe los tres corren sobre la real:

\begin{table}[H]
\centering
\caption{Las tres familias conviviendo en un solo catálogo}
\begin{tabularx}{\textwidth}{@{}llX@{}}
\toprule
\textbf{Familia} & \textbf{Dónde corre} & \textbf{Estructura que la atiende} \\
\midrule
\texttt{RTREE} & en el proceso del motor & R-Tree paginado en disco \\
\texttt{INVERTED} & en el proceso del motor & índice invertido en disco (SPIMI) \\
\texttt{IVF} & en un servicio aparte, por HTTP & IVF del bloque vectorial \\
\bottomrule
\end{tabularx}
\end{table}

\subsubsection{El índice que vive en otro proceso}

El IVF es el único que no corre dentro del motor. Sus vectores los produce un
modelo que arrastra opencv y onnxruntime, y meter eso en la imagen del motor la
multiplicaba de tamaño. El adaptador implementa el mismo contrato que los demás
y traduce por HTTP, con una complicación propia: el servicio identifica sus
rostros con un número correlativo y el motor trabaja con RIDs. El adaptador
construye el mapa entre los dos al crear el índice, leyendo la columna de rutas,
y es lo único que necesita para traducir en ambas direcciones.

Ese mapa es también lo único que se pierde al reiniciar, y encontrarlo nos costó.
Con el catálogo persistido, el motor arrancaba con el índice IVF declarado pero
vacío, así que la consulta vectorial se planificaba bien y devolvía cero filas.
Los índices espacial y textual tenían su rutina de reapertura y el vectorial no.
Ahora el arranque lo reconstruye recorriendo la tabla, que es donde las rutas ya
están guardadas, y hay una prueba que levanta el motor dos veces sobre el mismo
catálogo para fijar ese comportamiento.

\subsection{Planificador Híbrido}

\subsubsection{Las dos rutas}

Cuando una consulta trae un filtro relacional y además un ranking vectorial, hay
dos formas de resolverla y dan el mismo resultado por caminos de costo muy
distinto.

\begin{table}[H]
\centering
\caption{Las dos estrategias de una consulta híbrida}
\begin{tabularx}{\textwidth}{@{}llX@{}}
\toprule
\textbf{Estrategia} & \textbf{Nodo} & \textbf{Qué hace} \\
\midrule
Pre-filtering & \texttt{IVFRefine} & la vía relacional acota las filas y el índice
solo calcula la similitud de esas \\
Post-filtering & \texttt{IVFScan} con un \texttt{Filter} encima & el índice sondea
sus listas y devuelve candidatos ordenados; el filtro se aplica después \\
Ranking directo & \texttt{IVFScan} & no hay filtro relacional \\
\bottomrule
\end{tabularx}
\end{table}

La estrategia que corrió queda escrita en el nodo del plan y se ve tanto en la
consola SQL como en la pantalla de identificación, que es lo que el enunciado
pide reportar.

\subsubsection{La regla de decisión}

La primera versión comparaba cuántas filas deja el filtro contra cuántos
candidatos pediría el índice, y elegía pre-filtering cuando el filtro dejaba
menos. Suena razonable y está mal por un motivo que solo se ve midiendo: el
post-filtering aplica el filtro sobre lo que el índice devolvió, no sobre lo que
el índice miró. Si el filtro es muy selectivo, de los resultados del índice no
sobrevive casi ninguno y la consulta responde con menos filas de las pedidas sin
avisar de nada. Con \texttt{LIMIT 5} y un filtro que deja el 4 por ciento de la
tabla, nos devolvía una o dos.

La regla que quedó mira eso directamente. Si $C$ es el número de filas que el
índice puede entregar, $f$ la fracción de la tabla que pasa el filtro y $k$ las
filas pedidas, el post-filtering solo sirve si
\begin{equation}
    \left\lfloor C \cdot f \right\rfloor \ge m \cdot k ,
\end{equation}
con $m = 2$ de margen, porque cuántas sobreviven es un sorteo alrededor de
$C \cdot f$ y pedir lo justo falla la mitad de las veces. Cuando la condición no
se cumple, el motor acota antes aunque eso cueste recorrer la tabla entera.
Elegido el post-filtering, el nodo no pide $k$ filas al índice sino
\begin{equation}
    \min\left(C,\; m \cdot \left\lceil k / f \right\rceil\right),
\end{equation}
de manera que después del filtro sigan quedando las $k$. El plan deja ese número
a la vista en el campo \texttt{pide}.

Hay un tercer caso que al principio no habíamos previsto. Si el planificador no
puede estimar la selectividad del filtro, porque es una igualdad sobre una
columna sin estadísticas, no sabe si el post-filtering alcanza. Dar por bueno que
el filtro no reduce nada lleva a elegirlo y a devolver menos filas; por eso ante
una estimación que no aporta información el motor prefiere pre-filtering, que es
correcto siempre. La diferencia se nota en el filtro \texttt{lote = '00000'} de
nuestras pruebas, que deja mil filas y el catálogo no sabe cuántas.

\begin{lstlisting}[caption={Los dos planes para la misma foto de consulta}]
WHERE id < 100                          ->  pre-filtering
-> Project (columns=id, lote, similitud)
  -> Limit (rows=5)
    -> IVFRefine (index=ix_rostros_ivf dim=128 top_k=5)
      -> Filter (condition=id < 100)
        -> SeqScan (table=rostros)

WHERE id > 2                            ->  post-filtering
-> Project (columns=id, lote, similitud)
  -> Limit (rows=5)
    -> Filter (condition=id > 2)
      -> IVFScan (index=ix_rostros_ivf dim=128 top_k=5 pide=12 nprobe=4)
\end{lstlisting}

\subsubsection{Lo que solo apareció con el dataset completo}

Tres errores de esta parte sobrevivieron a las pruebas unitarias y aparecieron
recién al cargar los 13\,000 rostros. Los tres tenían la misma causa de fondo:
las pruebas corren sobre tablas de cinco filas, y a esa escala cualquier
constante parece razonable.

El primero fue que la decisión leía las filas que reportaba el plan relacional.
Un \texttt{SeqScan} con un \texttt{Filter} encima sigue informando que devuelve
la tabla entera, porque el filtro no cambia su estimación, así que el
planificador creía que ningún filtro reducía nada y elegía post-filtering
siempre. La selectividad hay que calcularla aparte, sobre los rangos sargables
del predicado.

El segundo fue la estimación de cuánto mira el índice. Estaba fijada en cuatro
veces $k$, que son veinte vectores para una consulta de cinco, cuando el IVF con
\texttt{nprobe = 4} sobre 128 listas sondea 1\,340 de 13\,000. Un orden de
magnitud de diferencia, y con el número chico el pre-filtering no ganaba nunca.
Ahora el servicio informa cuántas listas tiene su índice y la estimación sale de
ahí, de la fracción $\text{nprobe}/\text{listas}$, que no depende de $k$.

El tercero fue el tope de resultados. El servicio vectorial no devuelve más de
50 filas por consulta, un límite sensato para una interfaz de búsqueda y un techo
duro para el post-filtering. Pedir 326 candidatos no servía de nada y seguíamos
recibiendo dos filas. El planificador ahora conoce ese tope y lo usa como la $C$
de la fórmula, que es lo que movió la frontera a donde las mediciones la
encuentran.

\subsection{Aplicación de Identificación Facial}

La aplicación recibe una foto y responde con los rostros más parecidos del
conjunto indexado, pero no va al índice por su cuenta: construye una sentencia
SQL con el vector como literal y la ejecuta por el mismo camino que cualquier
otra consulta. Esa decisión tiene una consecuencia que vale el esfuerzo: la
pantalla puede mostrar el plan que corrió, los bloques que leyó y la estrategia
que el planificador eligió, porque son los mismos datos que devuelve la consola
SQL.

\begin{table}[H]
\centering
\caption{Endpoints añadidos al motor}
\begin{tabularx}{\textwidth}{@{}llX@{}}
\toprule
\textbf{Método} & \textbf{Ruta} & \textbf{Función} \\
\midrule
\texttt{GET} & \texttt{/api/identificacion/estado} & Si el bloque vectorial
responde y cuántos rostros tiene indexados \\
\texttt{POST} & \texttt{/api/identificacion/buscar} & Recibe la foto, arma la
consulta y devuelve filas, plan, estrategia y métricas \\
\texttt{GET} & \texttt{/api/identificacion/imagen} & Sirve una imagen indexada
por su ruta \\
\bottomrule
\end{tabularx}
\end{table}

El último endpoint existe por una razón de navegador. La galería de resultados
necesita mostrar las imágenes, que viven en el servicio vectorial en otro
puerto, y pedirlas desde ahí las bloqueaba por origen cruzado. El motor las
reenvía, con lo que el cliente habla con un solo origen igual que en el primer
entregable.

La pantalla tiene tres paneles: la foto de consulta con el filtro relacional y el
número de resultados, la galería ordenada por similitud con el primer puesto
marcado, y el panel que explica cómo se resolvió, con la sentencia que se
ejecutó, el reparto de tiempos y la ruta de acceso. Cambiar el filtro y ver cómo
el sello de la estrategia pasa de una a otra es la demostración más directa de
que el planificador está decidiendo de verdad.

\begin{table}[H]
\centering
\caption{Reparto de una identificación completa, sin filtro relacional}
\begin{tabular}{@{}lrl@{}}
\toprule
\textbf{Etapa} & \textbf{ms} & \textbf{Dónde ocurre} \\
\midrule
Vector de la foto (YuNet + SFace) & 48{,}50 & servicio vectorial \\
Análisis léxico y sintáctico & 0{,}39 & motor \\
Planificación & 0{,}06 & motor \\
Ejecución & 4{,}97 & motor y servicio \\
\midrule
\textbf{Total} & \textbf{53{,}92} & \\
\bottomrule
\end{tabular}
\end{table}

Nueve de cada diez milisegundos se los lleva calcular el vector de la foto, que
es detección más una pasada de la red de reconocimiento sobre una imagen de
1024 por 1024. Todo lo que hace el motor, incluido parsear una sentencia que
lleva 128 flotantes escritos como literal, suma cinco milisegundos. La
planificación cuesta 0{,}06, de modo que el trabajo que describe este informe no
es el que domina el tiempo de respuesta, y eso está bien: significa que la
consulta híbrida sale prácticamente gratis frente al costo de mirar la foto.

Sobre la calidad del resultado, la foto de consulta se recupera a sí misma con
similitud 1{,}000 y el segundo puesto queda en 0{,}373. Esa separación es la que
hace legible la galería: el primer rostro es el mismo y los siguientes son
parecidos de lejos.

\subsection{Experimento 9: Estrategia de Filtrado en Consultas Híbridas}

\begin{lstlisting}[language=bash, caption={Ejecución del experimento}]
python benchmarks/experimento9_hibrido.py
\end{lstlisting}

El arnés monta la tabla de 13\,000 rostros en un directorio temporal con el
backend de disco, toma una foto del propio conjunto como consulta y repite cada
medición cinco veces después de una corrida de calentamiento, quedándose con la
mediana. Deja tres archivos CSV en \texttt{benchmarks/results/}. Las cifras de
bloques son las que el \texttt{DiskCounter} registró.

\begin{table}[H]
\centering
\caption{Estrategia elegida según la selectividad del filtro, con $k = 5$}
\begin{tabular}{@{}lrrlrrr@{}}
\toprule
\textbf{Filtro} & \textbf{Filas} & \textbf{\% tabla} & \textbf{Estrategia} &
\textbf{Pide} & \textbf{Bloques} & \textbf{ms} \\
\midrule
sin filtro & 13\,000 & 100{,}00 & directo & 5 & 5 & 1{,}63 \\
\texttt{id \textless{} 10} & 10 & 0{,}08 & pre & 5 & 128 & 23{,}47 \\
\texttt{id \textless{} 100} & 100 & 0{,}77 & pre & 5 & 128 & 23{,}82 \\
\texttt{id \textless{} 500} & 500 & 3{,}85 & pre & 5 & 128 & 24{,}60 \\
\texttt{id \textless{} 1000} & 1\,000 & 7{,}69 & pre & 5 & 128 & 26{,}25 \\
\texttt{id \textless{} 2600} & 2\,600 & 20{,}00 & pre & 5 & 128 & 30{,}46 \\
\texttt{id \textless{} 5000} & 5\,000 & 38{,}46 & post & 28 & 28 & 1{,}79 \\
\texttt{id \textless{} 9000} & 9\,000 & 69{,}23 & post & 16 & 16 & 1{,}24 \\
\texttt{id \textgreater{} 2} & 12\,997 & 99{,}98 & post & 12 & 12 & 1{,}11 \\
\bottomrule
\end{tabular}
\end{table}

Las nueve consultas devolvieron las cinco filas pedidas, que es la invariante
que la regla de decisión y la sobre-emisión existen para sostener.

\begin{figure}[H]
\centering
\begin{tikzpicture}
\begin{axis}[figura,
    height=6.8cm,
    xmode=log, ymode=log,
    xlabel={Fracción de la tabla que deja pasar el filtro (\%)},
    ylabel={Bloques leídos},
    xmin=0.05, xmax=200, ymin=6, ymax=330,
    xtick={0.1,1,10,100},
    xticklabels={0{,}1\%,1\%,10\%,100\%},
    ytick={10,30,100,128},
    yticklabels={10,30,100,128},
    legend pos=south west,
]
% La referencia y la marca del cambio van antes que las series, para quedar
% por debajo de ellas.
\addplot[cHeap, dashed, mark=none, forget plot] coordinates {(0.05,128)(200,128)};
\node[nota, anchor=south west, text=cHeap] at (axis cs:0.055,133)
  {escaneo completo de la tabla};

\draw[cNota, dotted, line width=0.6pt] (axis cs:28,6) -- (axis cs:28,105);
\node[nota, anchor=east] at (axis cs:24,46)
  {el índice ya puede entregar\\las 5 filas tras el filtro};

\addplot[cBtree, mark=*]
  coordinates {(0.08,128)(0.77,128)(3.85,128)(7.69,128)(20,128)};
\addlegendentry{Pre-filtering (IVFRefine)}
\addplot[cSeq, mark=triangle*]
  coordinates {(38.46,28)(69.23,16)(99.98,12)};
\addlegendentry{Post-filtering (IVFScan)}
\end{axis}
\end{tikzpicture}
\caption{Bloques leídos según lo selectivo que sea el filtro relacional, con
$k = 5$ y ambos ejes logarítmicos. El pre-filtering paga el escaneo completo de
la tabla sin importar cuántas filas sobrevivan, porque \texttt{rostros} no tiene
índice sobre la columna del filtro; el post-filtering solo lee las filas que
termina devolviendo. La línea punteada vertical marca dónde el planificador
cambia de estrategia, que es donde el índice empieza a poder entregar las cinco
filas pedidas una vez aplicado el filtro.}
\label{fig:e9-estrategia}
\end{figure}

El salto entre las dos estrategias es de dos órdenes de magnitud en tiempo, de
30 milisegundos a 1{,}8, y de más de cuatro veces en bloques. La razón está en
que \texttt{rostros} no tiene índice sobre \texttt{id}: el pre-filtering acota
las filas recorriendo la tabla completa, así que su costo es el del escaneo
secuencial más el cálculo de similitud de las filas que pasaron. Es el precio de
poder responder la consulta completa, y el planificador lo paga solo cuando la
otra vía no puede.

\begin{table}[H]
\centering
\caption{Los dos extremos entre los que se mueven las estrategias}
\begin{tabular}{@{}lrr@{}}
\toprule
\textbf{Ruta} & \textbf{Bloques} & \textbf{ms} \\
\midrule
Escaneo completo con filtro (\texttt{SeqScan} + \texttt{Filter}) & 128 & 21{,}37 \\
Ranking por el IVF, sin filtro & 5 & 1{,}63 \\
\bottomrule
\end{tabular}
\end{table}

La tercera parte del experimento varía \texttt{nprobe}, que es cuántas de las
128 listas del índice se sondean. Lo interesante es que mueve mucho el trabajo
del índice y nada la estrategia:

\begin{table}[H]
\centering
\caption{Efecto de \texttt{nprobe} sobre el sondeo y sobre la decisión}
\begin{tabular}{@{}rrrrlr@{}}
\toprule
\textbf{nprobe} & \textbf{Candidatos} & \textbf{\% colección} &
\textbf{ms sin filtro} & \textbf{Estrategia con \texttt{id \textless{} 1000}} &
\textbf{Bloques} \\
\midrule
1 & 353 & 2{,}72 & 1{,}13 & pre & 128 \\
2 & 607 & 4{,}67 & 1{,}17 & pre & 128 \\
4 & 1\,340 & 10{,}31 & 1{,}22 & pre & 128 \\
8 & 2\,345 & 18{,}04 & 1{,}23 & pre & 128 \\
16 & 4\,474 & 34{,}42 & 1{,}59 & pre & 128 \\
\bottomrule
\end{tabular}
\end{table}

Los candidatos crecen de 353 a 4\,474, trece veces, y el tiempo sube de 1{,}13 a
1{,}59 milisegundos. La comparación de vectores está vectorizada en el servicio y
cuesta poco frente al resto, así que subir \texttt{nprobe} sale casi gratis en
tiempo aunque multiplique el trabajo nominal. La estrategia, en cambio, no se
mueve: con \texttt{id \textless{} 1000} el post-filtering seguiría sin alcanzar porque el
techo no es cuánto sondea el índice sino cuántas filas accede a devolver, que son
50. Sondear más no ayuda si el resultado se recorta igual.

Esa es la lección del experimento, y no la esperábamos. Pensábamos que la
decisión entre pre y post filtering iba a depender del costo de cada ruta, y
depende de algo anterior: de si la ruta puede responder la pregunta. El costo
solo decide entre las rutas que sí pueden, y en esta tabla, sin un índice sobre
la columna del filtro, cuando el post-filtering puede es además el más barato por
un margen amplio.

\subsection{Verificación}

La parte del motor tiene 248 pruebas automatizadas, de las cuales 12 cubren el
camino vectorial con un doble del servicio, de modo que corren sin opencv ni
modelos. Además hay una comprobación de extremo a extremo que levanta una tabla
por bloque, corre una consulta de cada familia y revisa que el plan elegido sea
el que corresponde:

\begin{lstlisting}[language=bash, caption={Comprobación de los cuatro bloques}]
python scripts/verificar_p2.py
\end{lstlisting}

\begin{lstlisting}[caption={Salida de la comprobación}]
Bloque A  espacial (R-Tree paginado)
  ok   busqueda por radio usa el R-Tree               3 locales
  ok   el radio deja fuera lo que no alcanza          ids [1, 2, 3]
  ok   k-NN usa el R-Tree                             2 < 3 < 1
Bloque B  textual (indice invertido SPIMI)
  ok   ranking por coseno usa el indice invertido     2 documentos
  ok   consulta booleana con NOT                      ids [5]
Bloque C  vectorial (IVF sobre rostros) y planificador hibrido
  ok   el motor ve el servicio vectorial              13000 imagenes, 12998 rostros
  ok   la foto de consulta se encuentra a si misma    00000/00000.png -> 100.0%
  ok   un filtro selectivo elige pre-filtering
  ok   un filtro debil elige post-filtering
Bloque D  catalogo y vertebra del motor
  ok   las tres familias de indice conviven en un catalogo  INVERTED IVF RTREE
  ok   los indices corren sobre su estructura real, no el sustituto
  ok   la via vectorial lee menos bloques que el escaneo completo  5 contra 129
\end{lstlisting}

El conjunto indexado son las carpetas \texttt{00000} a \texttt{13000} de FFHQ,
13\,000 imágenes de 1024 por 1024 que pesan 17 GB. El detector encontró rostro en
12\,998, y las dos que faltan son fotos donde la cara queda cortada o muy de
perfil.

%
% Si los experimentos se agrupan todos en la seccion de experimentacion, la
% subseccion "Experimento 9" se mueve alla sin tocar el resto.
% ==============================================================================

\section{Extensión Multimodal del Motor de Consultas}

En el primer entregable el motor resolvía un solo tipo de pregunta: comparar una
columna contra una constante. Este entregable le suma tres formas de buscar que
no se parecen entre sí ni a la anterior. La espacial pide los puntos cercanos a
una coordenada, la textual ordena documentos por relevancia frente a una
consulta de palabras, y la vectorial ordena imágenes por parecido frente a otra
imagen. Las tres devuelven filas ordenadas por una magnitud que la tabla no
guarda en ninguna columna.

El trabajo de esta parte fue que el motor aceptara las tres sin convertirse en
cuatro motores pegados. Lo que hicimos fue extender la gramática, agregar un
contrato común que cualquier estructura de índice puede implementar, y meter en
el planificador la decisión que el enunciado pide reportar: cuándo conviene
acotar las filas antes de consultar el índice y cuándo conviene lo contrario.
Sobre eso se construyó la aplicación de identificación facial, que es la que
pone los tres niveles a trabajar en una sola consulta.

\subsection{Gramática de las Consultas Multimodales}

\subsubsection{Las tres familias de predicado}

Las extensiones se añadieron al analizador léxico y al parser de descenso
recursivo que ya existían, sin cambiar la precedencia de operadores ni la forma
del árbol sintáctico. Cada familia entra por un sitio distinto de la gramática,
y esa diferencia no es estética: refleja dónde encaja cada una en la semántica
de un \texttt{SELECT}.

\begin{table}[H]
\centering
\caption{Dónde entra cada familia en la gramática}
\begin{tabularx}{\textwidth}{@{}llX@{}}
\toprule
\textbf{Familia} & \textbf{Lugar} & \textbf{Sintaxis} \\
\midrule
Espacial, radio & predicado del \texttt{WHERE} & \texttt{col WITHIN\_RADIUS((lon, lat), r)} \\
Espacial, k-NN & cláusula \texttt{ORDER BY} & \texttt{ORDER BY DISTANCE(col, (lon, lat))} \\
Textual & predicado del \texttt{WHERE} & \texttt{col MATCH 'términos'} \\
Textual, puntaje & lista de proyección & \texttt{COSINE\_SCORE() AS rel} \\
Vectorial & lista de proyección & \texttt{col \textless=\textgreater{} '[v1, ..., vn]' AS sim} \\
\bottomrule
\end{tabularx}
\end{table}

El radio es un predicado porque recorta el conjunto de filas y se puede combinar
con \texttt{AND} como cualquier otra condición. El k-NN, en cambio, no recorta
nada: ordena, y el recorte lo hace el \texttt{LIMIT}. Por eso vive en el
\texttt{ORDER BY} y no en el \texttt{WHERE}, que fue la primera cosa que tuvimos
que decidir y la que más ordenó el resto del diseño.

La distancia vectorial es el caso más incómodo de los cinco. No es un predicado
ni una cláusula de orden: es una columna calculada que no está en la tabla. El
operador \texttt{\textless=\textgreater{}} aparece en la proyección, produce una similitud por fila y
obliga a declarar un alias, porque sin nombre el \texttt{ORDER BY} no tiene cómo
referirse a ella. El motor rechaza la consulta si falta:

\begin{lstlisting}[language=SQL, caption={Declaración y uso del índice vectorial}]
CREATE TABLE rostros (
  id        INT PRIMARY KEY,
  lote      CHAR(8),
  embedding VARCHAR(23)
) USING HEAP;

CREATE INDEX ix_rostros_ivf ON rostros (embedding)
  USING IVF WITH (mode = 'face', nprobe = 4);

SELECT id, lote, embedding <=> '[0.111721, 0.152196, ...]' AS similitud
FROM rostros
WHERE id < 1000
ORDER BY similitud DESC
LIMIT 5;
\end{lstlisting}

\subsubsection{La columna indexada guarda una referencia, no el vector}

La tabla no almacena los 128 flotantes de cada rostro. Guarda la ruta de la
imagen, y es el índice el que resuelve de ahí el vector. Esto no fue una
concesión: es la misma idea con la que el R-Tree del bloque espacial indexa un
\texttt{POINT(lon lat)} escrito en una columna de texto. La tabla guarda la
referencia y la estructura guarda la geometría.

La consecuencia práctica es que un registro de \texttt{rostros} ocupa 36 bytes y
entran 101 por página, así que las 13\,000 filas caben en 129 páginas. Si el
vector viviera en la tabla, cada fila pesaría 512 bytes solo de flotantes y la
misma tabla ocuparía más de 1\,700 páginas, con el agravante de que ninguna
consulta relacional necesita leer esos bytes.

\subsection{Acoplamiento de las Tres Familias de Índice}

\subsubsection{Un contrato que no sabe de geometría ni de texto}

Las tres estructuras nuevas se enchufan por la misma interfaz que ya usaban el
árbol B$^+$ y el hash extensible. Un índice declara qué puede servir y el motor
le pregunta antes de registrarlo; si no puede, devuelve el motivo en texto y ese
motivo es lo que llega al usuario:

\begin{lstlisting}[caption={Motivo de rechazo de un índice}]
CREATE INDEX ix ON rostros (id) USING IVF;

  rostros.id debe guardar la ruta de la imagen en CHAR/VARCHAR
\end{lstlisting}

El despachador mantiene la lista de candidatos y encamina cada operación al que
aceptó ese índice. Lo que no acepta nadie cae a un sustituto en memoria, de modo
que una consulta sigue respondiendo aunque la estructura que le correspondía no
esté disponible. El cliente muestra en qué estructura quedó cada índice, y en
las mediciones de este informe los tres corren sobre la real:

\begin{table}[H]
\centering
\caption{Las tres familias conviviendo en un solo catálogo}
\begin{tabularx}{\textwidth}{@{}llX@{}}
\toprule
\textbf{Familia} & \textbf{Dónde corre} & \textbf{Estructura que la atiende} \\
\midrule
\texttt{RTREE} & en el proceso del motor & R-Tree paginado en disco \\
\texttt{INVERTED} & en el proceso del motor & índice invertido en disco (SPIMI) \\
\texttt{IVF} & en un servicio aparte, por HTTP & IVF del bloque vectorial \\
\bottomrule
\end{tabularx}
\end{table}

\subsubsection{El índice que vive en otro proceso}

El IVF es el único que no corre dentro del motor. Sus vectores los produce un
modelo que arrastra opencv y onnxruntime, y meter eso en la imagen del motor la
multiplicaba de tamaño. El adaptador implementa el mismo contrato que los demás
y traduce por HTTP, con una complicación propia: el servicio identifica sus
rostros con un número correlativo y el motor trabaja con RIDs. El adaptador
construye el mapa entre los dos al crear el índice, leyendo la columna de rutas,
y es lo único que necesita para traducir en ambas direcciones.

Ese mapa es también lo único que se pierde al reiniciar, y encontrarlo nos costó.
Con el catálogo persistido, el motor arrancaba con el índice IVF declarado pero
vacío, así que la consulta vectorial se planificaba bien y devolvía cero filas.
Los índices espacial y textual tenían su rutina de reapertura y el vectorial no.
Ahora el arranque lo reconstruye recorriendo la tabla, que es donde las rutas ya
están guardadas, y hay una prueba que levanta el motor dos veces sobre el mismo
catálogo para fijar ese comportamiento.

\subsection{Planificador Híbrido}

\subsubsection{Las dos rutas}

Cuando una consulta trae un filtro relacional y además un ranking vectorial, hay
dos formas de resolverla y dan el mismo resultado por caminos de costo muy
distinto.

\begin{table}[H]
\centering
\caption{Las dos estrategias de una consulta híbrida}
\begin{tabularx}{\textwidth}{@{}llX@{}}
\toprule
\textbf{Estrategia} & \textbf{Nodo} & \textbf{Qué hace} \\
\midrule
Pre-filtering & \texttt{IVFRefine} & la vía relacional acota las filas y el índice
solo calcula la similitud de esas \\
Post-filtering & \texttt{IVFScan} con un \texttt{Filter} encima & el índice sondea
sus listas y devuelve candidatos ordenados; el filtro se aplica después \\
Ranking directo & \texttt{IVFScan} & no hay filtro relacional \\
\bottomrule
\end{tabularx}
\end{table}

La estrategia que corrió queda escrita en el nodo del plan y se ve tanto en la
consola SQL como en la pantalla de identificación, que es lo que el enunciado
pide reportar.

\subsubsection{La regla de decisión}

La primera versión comparaba cuántas filas deja el filtro contra cuántos
candidatos pediría el índice, y elegía pre-filtering cuando el filtro dejaba
menos. Suena razonable y está mal por un motivo que solo se ve midiendo: el
post-filtering aplica el filtro sobre lo que el índice devolvió, no sobre lo que
el índice miró. Si el filtro es muy selectivo, de los resultados del índice no
sobrevive casi ninguno y la consulta responde con menos filas de las pedidas sin
avisar de nada. Con \texttt{LIMIT 5} y un filtro que deja el 4 por ciento de la
tabla, nos devolvía una o dos.

La regla que quedó mira eso directamente. Si $C$ es el número de filas que el
índice puede entregar, $f$ la fracción de la tabla que pasa el filtro y $k$ las
filas pedidas, el post-filtering solo sirve si
\begin{equation}
    \left\lfloor C \cdot f \right\rfloor \ge m \cdot k ,
\end{equation}
con $m = 2$ de margen, porque cuántas sobreviven es un sorteo alrededor de
$C \cdot f$ y pedir lo justo falla la mitad de las veces. Cuando la condición no
se cumple, el motor acota antes aunque eso cueste recorrer la tabla entera.
Elegido el post-filtering, el nodo no pide $k$ filas al índice sino
\begin{equation}
    \min\left(C,\; m \cdot \left\lceil k / f \right\rceil\right),
\end{equation}
de manera que después del filtro sigan quedando las $k$. El plan deja ese número
a la vista en el campo \texttt{pide}.

Hay un tercer caso que al principio no habíamos previsto. Si el planificador no
puede estimar la selectividad del filtro, porque es una igualdad sobre una
columna sin estadísticas, no sabe si el post-filtering alcanza. Dar por bueno que
el filtro no reduce nada lleva a elegirlo y a devolver menos filas; por eso ante
una estimación que no aporta información el motor prefiere pre-filtering, que es
correcto siempre. La diferencia se nota en el filtro \texttt{lote = '00000'} de
nuestras pruebas, que deja mil filas y el catálogo no sabe cuántas.

\begin{lstlisting}[caption={Los dos planes para la misma foto de consulta}]
WHERE id < 100                          ->  pre-filtering
-> Project (columns=id, lote, similitud)
  -> Limit (rows=5)
    -> IVFRefine (index=ix_rostros_ivf dim=128 top_k=5)
      -> Filter (condition=id < 100)
        -> SeqScan (table=rostros)

WHERE id > 2                            ->  post-filtering
-> Project (columns=id, lote, similitud)
  -> Limit (rows=5)
    -> Filter (condition=id > 2)
      -> IVFScan (index=ix_rostros_ivf dim=128 top_k=5 pide=12 nprobe=4)
\end{lstlisting}

\subsubsection{Lo que solo apareció con el dataset completo}

Tres errores de esta parte sobrevivieron a las pruebas unitarias y aparecieron
recién al cargar los 13\,000 rostros. Los tres tenían la misma causa de fondo:
las pruebas corren sobre tablas de cinco filas, y a esa escala cualquier
constante parece razonable.

El primero fue que la decisión leía las filas que reportaba el plan relacional.
Un \texttt{SeqScan} con un \texttt{Filter} encima sigue informando que devuelve
la tabla entera, porque el filtro no cambia su estimación, así que el
planificador creía que ningún filtro reducía nada y elegía post-filtering
siempre. La selectividad hay que calcularla aparte, sobre los rangos sargables
del predicado.

El segundo fue la estimación de cuánto mira el índice. Estaba fijada en cuatro
veces $k$, que son veinte vectores para una consulta de cinco, cuando el IVF con
\texttt{nprobe = 4} sobre 128 listas sondea 1\,340 de 13\,000. Un orden de
magnitud de diferencia, y con el número chico el pre-filtering no ganaba nunca.
Ahora el servicio informa cuántas listas tiene su índice y la estimación sale de
ahí, de la fracción $\text{nprobe}/\text{listas}$, que no depende de $k$.

El tercero fue el tope de resultados. El servicio vectorial no devuelve más de
50 filas por consulta, un límite sensato para una interfaz de búsqueda y un techo
duro para el post-filtering. Pedir 326 candidatos no servía de nada y seguíamos
recibiendo dos filas. El planificador ahora conoce ese tope y lo usa como la $C$
de la fórmula, que es lo que movió la frontera a donde las mediciones la
encuentran.

\subsection{Aplicación de Identificación Facial}

La aplicación recibe una foto y responde con los rostros más parecidos del
conjunto indexado, pero no va al índice por su cuenta: construye una sentencia
SQL con el vector como literal y la ejecuta por el mismo camino que cualquier
otra consulta. Esa decisión tiene una consecuencia que vale el esfuerzo: la
pantalla puede mostrar el plan que corrió, los bloques que leyó y la estrategia
que el planificador eligió, porque son los mismos datos que devuelve la consola
SQL.

\begin{table}[H]
\centering
\caption{Endpoints añadidos al motor}
\begin{tabularx}{\textwidth}{@{}llX@{}}
\toprule
\textbf{Método} & \textbf{Ruta} & \textbf{Función} \\
\midrule
\texttt{GET} & \texttt{/api/identificacion/estado} & Si el bloque vectorial
responde y cuántos rostros tiene indexados \\
\texttt{POST} & \texttt{/api/identificacion/buscar} & Recibe la foto, arma la
consulta y devuelve filas, plan, estrategia y métricas \\
\texttt{GET} & \texttt{/api/identificacion/imagen} & Sirve una imagen indexada
por su ruta \\
\bottomrule
\end{tabularx}
\end{table}

El último endpoint existe por una razón de navegador. La galería de resultados
necesita mostrar las imágenes, que viven en el servicio vectorial en otro
puerto, y pedirlas desde ahí las bloqueaba por origen cruzado. El motor las
reenvía, con lo que el cliente habla con un solo origen igual que en el primer
entregable.

La pantalla tiene tres paneles: la foto de consulta con el filtro relacional y el
número de resultados, la galería ordenada por similitud con el primer puesto
marcado, y el panel que explica cómo se resolvió, con la sentencia que se
ejecutó, el reparto de tiempos y la ruta de acceso. Cambiar el filtro y ver cómo
el sello de la estrategia pasa de una a otra es la demostración más directa de
que el planificador está decidiendo de verdad.

\begin{table}[H]
\centering
\caption{Reparto de una identificación completa, sin filtro relacional}
\begin{tabular}{@{}lrl@{}}
\toprule
\textbf{Etapa} & \textbf{ms} & \textbf{Dónde ocurre} \\
\midrule
Vector de la foto (YuNet + SFace) & 48{,}50 & servicio vectorial \\
Análisis léxico y sintáctico & 0{,}39 & motor \\
Planificación & 0{,}06 & motor \\
Ejecución & 4{,}97 & motor y servicio \\
\midrule
\textbf{Total} & \textbf{53{,}92} & \\
\bottomrule
\end{tabular}
\end{table}

Nueve de cada diez milisegundos se los lleva calcular el vector de la foto, que
es detección más una pasada de la red de reconocimiento sobre una imagen de
1024 por 1024. Todo lo que hace el motor, incluido parsear una sentencia que
lleva 128 flotantes escritos como literal, suma cinco milisegundos. La
planificación cuesta 0{,}06, de modo que el trabajo que describe este informe no
es el que domina el tiempo de respuesta, y eso está bien: significa que la
consulta híbrida sale prácticamente gratis frente al costo de mirar la foto.

Sobre la calidad del resultado, la foto de consulta se recupera a sí misma con
similitud 1{,}000 y el segundo puesto queda en 0{,}373. Esa separación es la que
hace legible la galería: el primer rostro es el mismo y los siguientes son
parecidos de lejos.

\subsection{Experimento 9: Estrategia de Filtrado en Consultas Híbridas}

\begin{lstlisting}[language=bash, caption={Ejecución del experimento}]
python benchmarks/experimento9_hibrido.py
\end{lstlisting}

El arnés monta la tabla de 13\,000 rostros en un directorio temporal con el
backend de disco, toma una foto del propio conjunto como consulta y repite cada
medición cinco veces después de una corrida de calentamiento, quedándose con la
mediana. Deja tres archivos CSV en \texttt{benchmarks/results/}. Las cifras de
bloques son las que el \texttt{DiskCounter} registró.

\begin{table}[H]
\centering
\caption{Estrategia elegida según la selectividad del filtro, con $k = 5$}
\begin{tabular}{@{}lrrlrrr@{}}
\toprule
\textbf{Filtro} & \textbf{Filas} & \textbf{\% tabla} & \textbf{Estrategia} &
\textbf{Pide} & \textbf{Bloques} & \textbf{ms} \\
\midrule
sin filtro & 13\,000 & 100{,}00 & directo & 5 & 5 & 1{,}63 \\
\texttt{id \textless{} 10} & 10 & 0{,}08 & pre & 5 & 128 & 23{,}47 \\
\texttt{id \textless{} 100} & 100 & 0{,}77 & pre & 5 & 128 & 23{,}82 \\
\texttt{id \textless{} 500} & 500 & 3{,}85 & pre & 5 & 128 & 24{,}60 \\
\texttt{id \textless{} 1000} & 1\,000 & 7{,}69 & pre & 5 & 128 & 26{,}25 \\
\texttt{id \textless{} 2600} & 2\,600 & 20{,}00 & pre & 5 & 128 & 30{,}46 \\
\texttt{id \textless{} 5000} & 5\,000 & 38{,}46 & post & 28 & 28 & 1{,}79 \\
\texttt{id \textless{} 9000} & 9\,000 & 69{,}23 & post & 16 & 16 & 1{,}24 \\
\texttt{id \textgreater{} 2} & 12\,997 & 99{,}98 & post & 12 & 12 & 1{,}11 \\
\bottomrule
\end{tabular}
\end{table}

Las nueve consultas devolvieron las cinco filas pedidas, que es la invariante
que la regla de decisión y la sobre-emisión existen para sostener.

\begin{figure}[H]
\centering
\begin{tikzpicture}
\begin{axis}[figura,
    height=6.8cm,
    xmode=log, ymode=log,
    xlabel={Fracción de la tabla que deja pasar el filtro (\%)},
    ylabel={Bloques leídos},
    xmin=0.05, xmax=200, ymin=6, ymax=330,
    xtick={0.1,1,10,100},
    xticklabels={0{,}1\%,1\%,10\%,100\%},
    ytick={10,30,100,128},
    yticklabels={10,30,100,128},
    legend pos=south west,
]
% La referencia y la marca del cambio van antes que las series, para quedar
% por debajo de ellas.
\addplot[cHeap, dashed, mark=none, forget plot] coordinates {(0.05,128)(200,128)};
\node[nota, anchor=south west, text=cHeap] at (axis cs:0.055,133)
  {escaneo completo de la tabla};

\draw[cNota, dotted, line width=0.6pt] (axis cs:28,6) -- (axis cs:28,105);
\node[nota, anchor=east] at (axis cs:24,46)
  {el índice ya puede entregar\\las 5 filas tras el filtro};

\addplot[cBtree, mark=*]
  coordinates {(0.08,128)(0.77,128)(3.85,128)(7.69,128)(20,128)};
\addlegendentry{Pre-filtering (IVFRefine)}
\addplot[cSeq, mark=triangle*]
  coordinates {(38.46,28)(69.23,16)(99.98,12)};
\addlegendentry{Post-filtering (IVFScan)}
\end{axis}
\end{tikzpicture}
\caption{Bloques leídos según lo selectivo que sea el filtro relacional, con
$k = 5$ y ambos ejes logarítmicos. El pre-filtering paga el escaneo completo de
la tabla sin importar cuántas filas sobrevivan, porque \texttt{rostros} no tiene
índice sobre la columna del filtro; el post-filtering solo lee las filas que
termina devolviendo. La línea punteada vertical marca dónde el planificador
cambia de estrategia, que es donde el índice empieza a poder entregar las cinco
filas pedidas una vez aplicado el filtro.}
\label{fig:e9-estrategia}
\end{figure}

El salto entre las dos estrategias es de dos órdenes de magnitud en tiempo, de
30 milisegundos a 1{,}8, y de más de cuatro veces en bloques. La razón está en
que \texttt{rostros} no tiene índice sobre \texttt{id}: el pre-filtering acota
las filas recorriendo la tabla completa, así que su costo es el del escaneo
secuencial más el cálculo de similitud de las filas que pasaron. Es el precio de
poder responder la consulta completa, y el planificador lo paga solo cuando la
otra vía no puede.

\begin{table}[H]
\centering
\caption{Los dos extremos entre los que se mueven las estrategias}
\begin{tabular}{@{}lrr@{}}
\toprule
\textbf{Ruta} & \textbf{Bloques} & \textbf{ms} \\
\midrule
Escaneo completo con filtro (\texttt{SeqScan} + \texttt{Filter}) & 128 & 21{,}37 \\
Ranking por el IVF, sin filtro & 5 & 1{,}63 \\
\bottomrule
\end{tabular}
\end{table}

La tercera parte del experimento varía \texttt{nprobe}, que es cuántas de las
128 listas del índice se sondean. Lo interesante es que mueve mucho el trabajo
del índice y nada la estrategia:

\begin{table}[H]
\centering
\caption{Efecto de \texttt{nprobe} sobre el sondeo y sobre la decisión}
\begin{tabular}{@{}rrrrlr@{}}
\toprule
\textbf{nprobe} & \textbf{Candidatos} & \textbf{\% colección} &
\textbf{ms sin filtro} & \textbf{Estrategia con \texttt{id \textless{} 1000}} &
\textbf{Bloques} \\
\midrule
1 & 353 & 2{,}72 & 1{,}13 & pre & 128 \\
2 & 607 & 4{,}67 & 1{,}17 & pre & 128 \\
4 & 1\,340 & 10{,}31 & 1{,}22 & pre & 128 \\
8 & 2\,345 & 18{,}04 & 1{,}23 & pre & 128 \\
16 & 4\,474 & 34{,}42 & 1{,}59 & pre & 128 \\
\bottomrule
\end{tabular}
\end{table}

Los candidatos crecen de 353 a 4\,474, trece veces, y el tiempo sube de 1{,}13 a
1{,}59 milisegundos. La comparación de vectores está vectorizada en el servicio y
cuesta poco frente al resto, así que subir \texttt{nprobe} sale casi gratis en
tiempo aunque multiplique el trabajo nominal. La estrategia, en cambio, no se
mueve: con \texttt{id \textless{} 1000} el post-filtering seguiría sin alcanzar porque el
techo no es cuánto sondea el índice sino cuántas filas accede a devolver, que son
50. Sondear más no ayuda si el resultado se recorta igual.

Esa es la lección del experimento, y no la esperábamos. Pensábamos que la
decisión entre pre y post filtering iba a depender del costo de cada ruta, y
depende de algo anterior: de si la ruta puede responder la pregunta. El costo
solo decide entre las rutas que sí pueden, y en esta tabla, sin un índice sobre
la columna del filtro, cuando el post-filtering puede es además el más barato por
un margen amplio.

\subsection{Verificación}

La parte del motor tiene 248 pruebas automatizadas, de las cuales 12 cubren el
camino vectorial con un doble del servicio, de modo que corren sin opencv ni
modelos. Además hay una comprobación de extremo a extremo que levanta una tabla
por bloque, corre una consulta de cada familia y revisa que el plan elegido sea
el que corresponde:

\begin{lstlisting}[language=bash, caption={Comprobación de los cuatro bloques}]
python scripts/verificar_p2.py
\end{lstlisting}

\begin{lstlisting}[caption={Salida de la comprobación}]
Bloque A  espacial (R-Tree paginado)
  ok   busqueda por radio usa el R-Tree               3 locales
  ok   el radio deja fuera lo que no alcanza          ids [1, 2, 3]
  ok   k-NN usa el R-Tree                             2 < 3 < 1
Bloque B  textual (indice invertido SPIMI)
  ok   ranking por coseno usa el indice invertido     2 documentos
  ok   consulta booleana con NOT                      ids [5]
Bloque C  vectorial (IVF sobre rostros) y planificador hibrido
  ok   el motor ve el servicio vectorial              13000 imagenes, 12998 rostros
  ok   la foto de consulta se encuentra a si misma    00000/00000.png -> 100.0%
  ok   un filtro selectivo elige pre-filtering
  ok   un filtro debil elige post-filtering
Bloque D  catalogo y vertebra del motor
  ok   las tres familias de indice conviven en un catalogo  INVERTED IVF RTREE
  ok   los indices corren sobre su estructura real, no el sustituto
  ok   la via vectorial lee menos bloques que el escaneo completo  5 contra 129
\end{lstlisting}

El conjunto indexado son las carpetas \texttt{00000} a \texttt{13000} de FFHQ,
13\,000 imágenes de 1024 por 1024 que pesan 17 GB. El detector encontró rostro en
12\,998, y las dos que faltan son fotos donde la cara queda cortada o muy de
perfil.

%
% Si los experimentos se agrupan todos en la seccion de experimentacion, la
% subseccion "Experimento 9" se mueve alla sin tocar el resto.
% ==============================================================================

\section{Extensión Multimodal del Motor de Consultas}

En el primer entregable el motor resolvía un solo tipo de pregunta: comparar una
columna contra una constante. Este entregable le suma tres formas de buscar que
no se parecen entre sí ni a la anterior. La espacial pide los puntos cercanos a
una coordenada, la textual ordena documentos por relevancia frente a una
consulta de palabras, y la vectorial ordena imágenes por parecido frente a otra
imagen. Las tres devuelven filas ordenadas por una magnitud que la tabla no
guarda en ninguna columna.

El trabajo de esta parte fue que el motor aceptara las tres sin convertirse en
cuatro motores pegados. Lo que hicimos fue extender la gramática, agregar un
contrato común que cualquier estructura de índice puede implementar, y meter en
el planificador la decisión que el enunciado pide reportar: cuándo conviene
acotar las filas antes de consultar el índice y cuándo conviene lo contrario.
Sobre eso se construyó la aplicación de identificación facial, que es la que
pone los tres niveles a trabajar en una sola consulta.

\subsection{Gramática de las Consultas Multimodales}

\subsubsection{Las tres familias de predicado}

Las extensiones se añadieron al analizador léxico y al parser de descenso
recursivo que ya existían, sin cambiar la precedencia de operadores ni la forma
del árbol sintáctico. Cada familia entra por un sitio distinto de la gramática,
y esa diferencia no es estética: refleja dónde encaja cada una en la semántica
de un \texttt{SELECT}.

\begin{table}[H]
\centering
\caption{Dónde entra cada familia en la gramática}
\begin{tabularx}{\textwidth}{@{}llX@{}}
\toprule
\textbf{Familia} & \textbf{Lugar} & \textbf{Sintaxis} \\
\midrule
Espacial, radio & predicado del \texttt{WHERE} & \texttt{col WITHIN\_RADIUS((lon, lat), r)} \\
Espacial, k-NN & cláusula \texttt{ORDER BY} & \texttt{ORDER BY DISTANCE(col, (lon, lat))} \\
Textual & predicado del \texttt{WHERE} & \texttt{col MATCH 'términos'} \\
Textual, puntaje & lista de proyección & \texttt{COSINE\_SCORE() AS rel} \\
Vectorial & lista de proyección & \texttt{col \textless=\textgreater{} '[v1, ..., vn]' AS sim} \\
\bottomrule
\end{tabularx}
\end{table}

El radio es un predicado porque recorta el conjunto de filas y se puede combinar
con \texttt{AND} como cualquier otra condición. El k-NN, en cambio, no recorta
nada: ordena, y el recorte lo hace el \texttt{LIMIT}. Por eso vive en el
\texttt{ORDER BY} y no en el \texttt{WHERE}, que fue la primera cosa que tuvimos
que decidir y la que más ordenó el resto del diseño.

La distancia vectorial es el caso más incómodo de los cinco. No es un predicado
ni una cláusula de orden: es una columna calculada que no está en la tabla. El
operador \texttt{\textless=\textgreater{}} aparece en la proyección, produce una similitud por fila y
obliga a declarar un alias, porque sin nombre el \texttt{ORDER BY} no tiene cómo
referirse a ella. El motor rechaza la consulta si falta:

\begin{lstlisting}[language=SQL, caption={Declaración y uso del índice vectorial}]
CREATE TABLE rostros (
  id        INT PRIMARY KEY,
  lote      CHAR(8),
  embedding VARCHAR(23)
) USING HEAP;

CREATE INDEX ix_rostros_ivf ON rostros (embedding)
  USING IVF WITH (mode = 'face', nprobe = 4);

SELECT id, lote, embedding <=> '[0.111721, 0.152196, ...]' AS similitud
FROM rostros
WHERE id < 1000
ORDER BY similitud DESC
LIMIT 5;
\end{lstlisting}

\subsubsection{La columna indexada guarda una referencia, no el vector}

La tabla no almacena los 128 flotantes de cada rostro. Guarda la ruta de la
imagen, y es el índice el que resuelve de ahí el vector. Esto no fue una
concesión: es la misma idea con la que el R-Tree del bloque espacial indexa un
\texttt{POINT(lon lat)} escrito en una columna de texto. La tabla guarda la
referencia y la estructura guarda la geometría.

La consecuencia práctica es que un registro de \texttt{rostros} ocupa 36 bytes y
entran 101 por página, así que las 13\,000 filas caben en 129 páginas. Si el
vector viviera en la tabla, cada fila pesaría 512 bytes solo de flotantes y la
misma tabla ocuparía más de 1\,700 páginas, con el agravante de que ninguna
consulta relacional necesita leer esos bytes.

\subsection{Acoplamiento de las Tres Familias de Índice}

\subsubsection{Un contrato que no sabe de geometría ni de texto}

Las tres estructuras nuevas se enchufan por la misma interfaz que ya usaban el
árbol B$^+$ y el hash extensible. Un índice declara qué puede servir y el motor
le pregunta antes de registrarlo; si no puede, devuelve el motivo en texto y ese
motivo es lo que llega al usuario:

\begin{lstlisting}[caption={Motivo de rechazo de un índice}]
CREATE INDEX ix ON rostros (id) USING IVF;

  rostros.id debe guardar la ruta de la imagen en CHAR/VARCHAR
\end{lstlisting}

El despachador mantiene la lista de candidatos y encamina cada operación al que
aceptó ese índice. Lo que no acepta nadie cae a un sustituto en memoria, de modo
que una consulta sigue respondiendo aunque la estructura que le correspondía no
esté disponible. El cliente muestra en qué estructura quedó cada índice, y en
las mediciones de este informe los tres corren sobre la real:

\begin{table}[H]
\centering
\caption{Las tres familias conviviendo en un solo catálogo}
\begin{tabularx}{\textwidth}{@{}llX@{}}
\toprule
\textbf{Familia} & \textbf{Dónde corre} & \textbf{Estructura que la atiende} \\
\midrule
\texttt{RTREE} & en el proceso del motor & R-Tree paginado en disco \\
\texttt{INVERTED} & en el proceso del motor & índice invertido en disco (SPIMI) \\
\texttt{IVF} & en un servicio aparte, por HTTP & IVF del bloque vectorial \\
\bottomrule
\end{tabularx}
\end{table}

\subsubsection{El índice que vive en otro proceso}

El IVF es el único que no corre dentro del motor. Sus vectores los produce un
modelo que arrastra opencv y onnxruntime, y meter eso en la imagen del motor la
multiplicaba de tamaño. El adaptador implementa el mismo contrato que los demás
y traduce por HTTP, con una complicación propia: el servicio identifica sus
rostros con un número correlativo y el motor trabaja con RIDs. El adaptador
construye el mapa entre los dos al crear el índice, leyendo la columna de rutas,
y es lo único que necesita para traducir en ambas direcciones.

Ese mapa es también lo único que se pierde al reiniciar, y encontrarlo nos costó.
Con el catálogo persistido, el motor arrancaba con el índice IVF declarado pero
vacío, así que la consulta vectorial se planificaba bien y devolvía cero filas.
Los índices espacial y textual tenían su rutina de reapertura y el vectorial no.
Ahora el arranque lo reconstruye recorriendo la tabla, que es donde las rutas ya
están guardadas, y hay una prueba que levanta el motor dos veces sobre el mismo
catálogo para fijar ese comportamiento.

\subsection{Planificador Híbrido}

\subsubsection{Las dos rutas}

Cuando una consulta trae un filtro relacional y además un ranking vectorial, hay
dos formas de resolverla y dan el mismo resultado por caminos de costo muy
distinto.

\begin{table}[H]
\centering
\caption{Las dos estrategias de una consulta híbrida}
\begin{tabularx}{\textwidth}{@{}llX@{}}
\toprule
\textbf{Estrategia} & \textbf{Nodo} & \textbf{Qué hace} \\
\midrule
Pre-filtering & \texttt{IVFRefine} & la vía relacional acota las filas y el índice
solo calcula la similitud de esas \\
Post-filtering & \texttt{IVFScan} con un \texttt{Filter} encima & el índice sondea
sus listas y devuelve candidatos ordenados; el filtro se aplica después \\
Ranking directo & \texttt{IVFScan} & no hay filtro relacional \\
\bottomrule
\end{tabularx}
\end{table}

La estrategia que corrió queda escrita en el nodo del plan y se ve tanto en la
consola SQL como en la pantalla de identificación, que es lo que el enunciado
pide reportar.

\subsubsection{La regla de decisión}

La primera versión comparaba cuántas filas deja el filtro contra cuántos
candidatos pediría el índice, y elegía pre-filtering cuando el filtro dejaba
menos. Suena razonable y está mal por un motivo que solo se ve midiendo: el
post-filtering aplica el filtro sobre lo que el índice devolvió, no sobre lo que
el índice miró. Si el filtro es muy selectivo, de los resultados del índice no
sobrevive casi ninguno y la consulta responde con menos filas de las pedidas sin
avisar de nada. Con \texttt{LIMIT 5} y un filtro que deja el 4 por ciento de la
tabla, nos devolvía una o dos.

La regla que quedó mira eso directamente. Si $C$ es el número de filas que el
índice puede entregar, $f$ la fracción de la tabla que pasa el filtro y $k$ las
filas pedidas, el post-filtering solo sirve si
\begin{equation}
    \left\lfloor C \cdot f \right\rfloor \ge m \cdot k ,
\end{equation}
con $m = 2$ de margen, porque cuántas sobreviven es un sorteo alrededor de
$C \cdot f$ y pedir lo justo falla la mitad de las veces. Cuando la condición no
se cumple, el motor acota antes aunque eso cueste recorrer la tabla entera.
Elegido el post-filtering, el nodo no pide $k$ filas al índice sino
\begin{equation}
    \min\left(C,\; m \cdot \left\lceil k / f \right\rceil\right),
\end{equation}
de manera que después del filtro sigan quedando las $k$. El plan deja ese número
a la vista en el campo \texttt{pide}.

Hay un tercer caso que al principio no habíamos previsto. Si el planificador no
puede estimar la selectividad del filtro, porque es una igualdad sobre una
columna sin estadísticas, no sabe si el post-filtering alcanza. Dar por bueno que
el filtro no reduce nada lleva a elegirlo y a devolver menos filas; por eso ante
una estimación que no aporta información el motor prefiere pre-filtering, que es
correcto siempre. La diferencia se nota en el filtro \texttt{lote = '00000'} de
nuestras pruebas, que deja mil filas y el catálogo no sabe cuántas.

\begin{lstlisting}[caption={Los dos planes para la misma foto de consulta}]
WHERE id < 100                          ->  pre-filtering
-> Project (columns=id, lote, similitud)
  -> Limit (rows=5)
    -> IVFRefine (index=ix_rostros_ivf dim=128 top_k=5)
      -> Filter (condition=id < 100)
        -> SeqScan (table=rostros)

WHERE id > 2                            ->  post-filtering
-> Project (columns=id, lote, similitud)
  -> Limit (rows=5)
    -> Filter (condition=id > 2)
      -> IVFScan (index=ix_rostros_ivf dim=128 top_k=5 pide=12 nprobe=4)
\end{lstlisting}

\subsubsection{Lo que solo apareció con el dataset completo}

Tres errores de esta parte sobrevivieron a las pruebas unitarias y aparecieron
recién al cargar los 13\,000 rostros. Los tres tenían la misma causa de fondo:
las pruebas corren sobre tablas de cinco filas, y a esa escala cualquier
constante parece razonable.

El primero fue que la decisión leía las filas que reportaba el plan relacional.
Un \texttt{SeqScan} con un \texttt{Filter} encima sigue informando que devuelve
la tabla entera, porque el filtro no cambia su estimación, así que el
planificador creía que ningún filtro reducía nada y elegía post-filtering
siempre. La selectividad hay que calcularla aparte, sobre los rangos sargables
del predicado.

El segundo fue la estimación de cuánto mira el índice. Estaba fijada en cuatro
veces $k$, que son veinte vectores para una consulta de cinco, cuando el IVF con
\texttt{nprobe = 4} sobre 128 listas sondea 1\,340 de 13\,000. Un orden de
magnitud de diferencia, y con el número chico el pre-filtering no ganaba nunca.
Ahora el servicio informa cuántas listas tiene su índice y la estimación sale de
ahí, de la fracción $\text{nprobe}/\text{listas}$, que no depende de $k$.

El tercero fue el tope de resultados. El servicio vectorial no devuelve más de
50 filas por consulta, un límite sensato para una interfaz de búsqueda y un techo
duro para el post-filtering. Pedir 326 candidatos no servía de nada y seguíamos
recibiendo dos filas. El planificador ahora conoce ese tope y lo usa como la $C$
de la fórmula, que es lo que movió la frontera a donde las mediciones la
encuentran.

\subsection{Aplicación de Identificación Facial}

La aplicación recibe una foto y responde con los rostros más parecidos del
conjunto indexado, pero no va al índice por su cuenta: construye una sentencia
SQL con el vector como literal y la ejecuta por el mismo camino que cualquier
otra consulta. Esa decisión tiene una consecuencia que vale el esfuerzo: la
pantalla puede mostrar el plan que corrió, los bloques que leyó y la estrategia
que el planificador eligió, porque son los mismos datos que devuelve la consola
SQL.

\begin{table}[H]
\centering
\caption{Endpoints añadidos al motor}
\begin{tabularx}{\textwidth}{@{}llX@{}}
\toprule
\textbf{Método} & \textbf{Ruta} & \textbf{Función} \\
\midrule
\texttt{GET} & \texttt{/api/identificacion/estado} & Si el bloque vectorial
responde y cuántos rostros tiene indexados \\
\texttt{POST} & \texttt{/api/identificacion/buscar} & Recibe la foto, arma la
consulta y devuelve filas, plan, estrategia y métricas \\
\texttt{GET} & \texttt{/api/identificacion/imagen} & Sirve una imagen indexada
por su ruta \\
\bottomrule
\end{tabularx}
\end{table}

El último endpoint existe por una razón de navegador. La galería de resultados
necesita mostrar las imágenes, que viven en el servicio vectorial en otro
puerto, y pedirlas desde ahí las bloqueaba por origen cruzado. El motor las
reenvía, con lo que el cliente habla con un solo origen igual que en el primer
entregable.

La pantalla tiene tres paneles: la foto de consulta con el filtro relacional y el
número de resultados, la galería ordenada por similitud con el primer puesto
marcado, y el panel que explica cómo se resolvió, con la sentencia que se
ejecutó, el reparto de tiempos y la ruta de acceso. Cambiar el filtro y ver cómo
el sello de la estrategia pasa de una a otra es la demostración más directa de
que el planificador está decidiendo de verdad.

\begin{table}[H]
\centering
\caption{Reparto de una identificación completa, sin filtro relacional}
\begin{tabular}{@{}lrl@{}}
\toprule
\textbf{Etapa} & \textbf{ms} & \textbf{Dónde ocurre} \\
\midrule
Vector de la foto (YuNet + SFace) & 48{,}50 & servicio vectorial \\
Análisis léxico y sintáctico & 0{,}39 & motor \\
Planificación & 0{,}06 & motor \\
Ejecución & 4{,}97 & motor y servicio \\
\midrule
\textbf{Total} & \textbf{53{,}92} & \\
\bottomrule
\end{tabular}
\end{table}

Nueve de cada diez milisegundos se los lleva calcular el vector de la foto, que
es detección más una pasada de la red de reconocimiento sobre una imagen de
1024 por 1024. Todo lo que hace el motor, incluido parsear una sentencia que
lleva 128 flotantes escritos como literal, suma cinco milisegundos. La
planificación cuesta 0{,}06, de modo que el trabajo que describe este informe no
es el que domina el tiempo de respuesta, y eso está bien: significa que la
consulta híbrida sale prácticamente gratis frente al costo de mirar la foto.

Sobre la calidad del resultado, la foto de consulta se recupera a sí misma con
similitud 1{,}000 y el segundo puesto queda en 0{,}373. Esa separación es la que
hace legible la galería: el primer rostro es el mismo y los siguientes son
parecidos de lejos.

\subsection{Experimento 9: Estrategia de Filtrado en Consultas Híbridas}

\begin{lstlisting}[language=bash, caption={Ejecución del experimento}]
python benchmarks/experimento9_hibrido.py
\end{lstlisting}

El arnés monta la tabla de 13\,000 rostros en un directorio temporal con el
backend de disco, toma una foto del propio conjunto como consulta y repite cada
medición cinco veces después de una corrida de calentamiento, quedándose con la
mediana. Deja tres archivos CSV en \texttt{benchmarks/results/}. Las cifras de
bloques son las que el \texttt{DiskCounter} registró.

\begin{table}[H]
\centering
\caption{Estrategia elegida según la selectividad del filtro, con $k = 5$}
\begin{tabular}{@{}lrrlrrr@{}}
\toprule
\textbf{Filtro} & \textbf{Filas} & \textbf{\% tabla} & \textbf{Estrategia} &
\textbf{Pide} & \textbf{Bloques} & \textbf{ms} \\
\midrule
sin filtro & 13\,000 & 100{,}00 & directo & 5 & 5 & 1{,}63 \\
\texttt{id \textless{} 10} & 10 & 0{,}08 & pre & 5 & 128 & 23{,}47 \\
\texttt{id \textless{} 100} & 100 & 0{,}77 & pre & 5 & 128 & 23{,}82 \\
\texttt{id \textless{} 500} & 500 & 3{,}85 & pre & 5 & 128 & 24{,}60 \\
\texttt{id \textless{} 1000} & 1\,000 & 7{,}69 & pre & 5 & 128 & 26{,}25 \\
\texttt{id \textless{} 2600} & 2\,600 & 20{,}00 & pre & 5 & 128 & 30{,}46 \\
\texttt{id \textless{} 5000} & 5\,000 & 38{,}46 & post & 28 & 28 & 1{,}79 \\
\texttt{id \textless{} 9000} & 9\,000 & 69{,}23 & post & 16 & 16 & 1{,}24 \\
\texttt{id \textgreater{} 2} & 12\,997 & 99{,}98 & post & 12 & 12 & 1{,}11 \\
\bottomrule
\end{tabular}
\end{table}

Las nueve consultas devolvieron las cinco filas pedidas, que es la invariante
que la regla de decisión y la sobre-emisión existen para sostener.

\begin{figure}[H]
\centering
\begin{tikzpicture}
\begin{axis}[figura,
    height=6.8cm,
    xmode=log, ymode=log,
    xlabel={Fracción de la tabla que deja pasar el filtro (\%)},
    ylabel={Bloques leídos},
    xmin=0.05, xmax=200, ymin=6, ymax=330,
    xtick={0.1,1,10,100},
    xticklabels={0{,}1\%,1\%,10\%,100\%},
    ytick={10,30,100,128},
    yticklabels={10,30,100,128},
    legend pos=south west,
]
% La referencia y la marca del cambio van antes que las series, para quedar
% por debajo de ellas.
\addplot[cHeap, dashed, mark=none, forget plot] coordinates {(0.05,128)(200,128)};
\node[nota, anchor=south west, text=cHeap] at (axis cs:0.055,133)
  {escaneo completo de la tabla};

\draw[cNota, dotted, line width=0.6pt] (axis cs:28,6) -- (axis cs:28,105);
\node[nota, anchor=east] at (axis cs:24,46)
  {el índice ya puede entregar\\las 5 filas tras el filtro};

\addplot[cBtree, mark=*]
  coordinates {(0.08,128)(0.77,128)(3.85,128)(7.69,128)(20,128)};
\addlegendentry{Pre-filtering (IVFRefine)}
\addplot[cSeq, mark=triangle*]
  coordinates {(38.46,28)(69.23,16)(99.98,12)};
\addlegendentry{Post-filtering (IVFScan)}
\end{axis}
\end{tikzpicture}
\caption{Bloques leídos según lo selectivo que sea el filtro relacional, con
$k = 5$ y ambos ejes logarítmicos. El pre-filtering paga el escaneo completo de
la tabla sin importar cuántas filas sobrevivan, porque \texttt{rostros} no tiene
índice sobre la columna del filtro; el post-filtering solo lee las filas que
termina devolviendo. La línea punteada vertical marca dónde el planificador
cambia de estrategia, que es donde el índice empieza a poder entregar las cinco
filas pedidas una vez aplicado el filtro.}
\label{fig:e9-estrategia}
\end{figure}

El salto entre las dos estrategias es de dos órdenes de magnitud en tiempo, de
30 milisegundos a 1{,}8, y de más de cuatro veces en bloques. La razón está en
que \texttt{rostros} no tiene índice sobre \texttt{id}: el pre-filtering acota
las filas recorriendo la tabla completa, así que su costo es el del escaneo
secuencial más el cálculo de similitud de las filas que pasaron. Es el precio de
poder responder la consulta completa, y el planificador lo paga solo cuando la
otra vía no puede.

\begin{table}[H]
\centering
\caption{Los dos extremos entre los que se mueven las estrategias}
\begin{tabular}{@{}lrr@{}}
\toprule
\textbf{Ruta} & \textbf{Bloques} & \textbf{ms} \\
\midrule
Escaneo completo con filtro (\texttt{SeqScan} + \texttt{Filter}) & 128 & 21{,}37 \\
Ranking por el IVF, sin filtro & 5 & 1{,}63 \\
\bottomrule
\end{tabular}
\end{table}

La tercera parte del experimento varía \texttt{nprobe}, que es cuántas de las
128 listas del índice se sondean. Lo interesante es que mueve mucho el trabajo
del índice y nada la estrategia:

\begin{table}[H]
\centering
\caption{Efecto de \texttt{nprobe} sobre el sondeo y sobre la decisión}
\begin{tabular}{@{}rrrrlr@{}}
\toprule
\textbf{nprobe} & \textbf{Candidatos} & \textbf{\% colección} &
\textbf{ms sin filtro} & \textbf{Estrategia con \texttt{id \textless{} 1000}} &
\textbf{Bloques} \\
\midrule
1 & 353 & 2{,}72 & 1{,}13 & pre & 128 \\
2 & 607 & 4{,}67 & 1{,}17 & pre & 128 \\
4 & 1\,340 & 10{,}31 & 1{,}22 & pre & 128 \\
8 & 2\,345 & 18{,}04 & 1{,}23 & pre & 128 \\
16 & 4\,474 & 34{,}42 & 1{,}59 & pre & 128 \\
\bottomrule
\end{tabular}
\end{table}

Los candidatos crecen de 353 a 4\,474, trece veces, y el tiempo sube de 1{,}13 a
1{,}59 milisegundos. La comparación de vectores está vectorizada en el servicio y
cuesta poco frente al resto, así que subir \texttt{nprobe} sale casi gratis en
tiempo aunque multiplique el trabajo nominal. La estrategia, en cambio, no se
mueve: con \texttt{id \textless{} 1000} el post-filtering seguiría sin alcanzar porque el
techo no es cuánto sondea el índice sino cuántas filas accede a devolver, que son
50. Sondear más no ayuda si el resultado se recorta igual.

Esa es la lección del experimento, y no la esperábamos. Pensábamos que la
decisión entre pre y post filtering iba a depender del costo de cada ruta, y
depende de algo anterior: de si la ruta puede responder la pregunta. El costo
solo decide entre las rutas que sí pueden, y en esta tabla, sin un índice sobre
la columna del filtro, cuando el post-filtering puede es además el más barato por
un margen amplio.

\subsection{Verificación}

La parte del motor tiene 248 pruebas automatizadas, de las cuales 12 cubren el
camino vectorial con un doble del servicio, de modo que corren sin opencv ni
modelos. Además hay una comprobación de extremo a extremo que levanta una tabla
por bloque, corre una consulta de cada familia y revisa que el plan elegido sea
el que corresponde:

\begin{lstlisting}[language=bash, caption={Comprobación de los cuatro bloques}]
python scripts/verificar_p2.py
\end{lstlisting}

\begin{lstlisting}[caption={Salida de la comprobación}]
Bloque A  espacial (R-Tree paginado)
  ok   busqueda por radio usa el R-Tree               3 locales
  ok   el radio deja fuera lo que no alcanza          ids [1, 2, 3]
  ok   k-NN usa el R-Tree                             2 < 3 < 1
Bloque B  textual (indice invertido SPIMI)
  ok   ranking por coseno usa el indice invertido     2 documentos
  ok   consulta booleana con NOT                      ids [5]
Bloque C  vectorial (IVF sobre rostros) y planificador hibrido
  ok   el motor ve el servicio vectorial              13000 imagenes, 12998 rostros
  ok   la foto de consulta se encuentra a si misma    00000/00000.png -> 100.0%
  ok   un filtro selectivo elige pre-filtering
  ok   un filtro debil elige post-filtering
Bloque D  catalogo y vertebra del motor
  ok   las tres familias de indice conviven en un catalogo  INVERTED IVF RTREE
  ok   los indices corren sobre su estructura real, no el sustituto
  ok   la via vectorial lee menos bloques que el escaneo completo  5 contra 129
\end{lstlisting}

El conjunto indexado son las carpetas \texttt{00000} a \texttt{13000} de FFHQ,
13\,000 imágenes de 1024 por 1024 que pesan 17 GB. El detector encontró rostro en
12\,998, y las dos que faltan son fotos donde la cara queda cortada o muy de
perfil.

%
% Si los experimentos se agrupan todos en la seccion de experimentacion, la
% subseccion "Experimento 9" se mueve alla sin tocar el resto.
% ==============================================================================

\section{Extensión Multimodal del Motor de Consultas}

En el primer entregable el motor resolvía un solo tipo de pregunta: comparar una
columna contra una constante. Este entregable le suma tres formas de buscar que
no se parecen entre sí ni a la anterior. La espacial pide los puntos cercanos a
una coordenada, la textual ordena documentos por relevancia frente a una
consulta de palabras, y la vectorial ordena imágenes por parecido frente a otra
imagen. Las tres devuelven filas ordenadas por una magnitud que la tabla no
guarda en ninguna columna.

El trabajo de esta parte fue que el motor aceptara las tres sin convertirse en
cuatro motores pegados. Lo que hicimos fue extender la gramática, agregar un
contrato común que cualquier estructura de índice puede implementar, y meter en
el planificador la decisión que el enunciado pide reportar: cuándo conviene
acotar las filas antes de consultar el índice y cuándo conviene lo contrario.
Sobre eso se construyó la aplicación de identificación facial, que es la que
pone los tres niveles a trabajar en una sola consulta.

\subsection{Gramática de las Consultas Multimodales}

\subsubsection{Las tres familias de predicado}

Las extensiones se añadieron al analizador léxico y al parser de descenso
recursivo que ya existían, sin cambiar la precedencia de operadores ni la forma
del árbol sintáctico. Cada familia entra por un sitio distinto de la gramática,
y esa diferencia no es estética: refleja dónde encaja cada una en la semántica
de un \texttt{SELECT}.

\begin{table}[H]
\centering
\caption{Dónde entra cada familia en la gramática}
\begin{tabularx}{\textwidth}{@{}llX@{}}
\toprule
\textbf{Familia} & \textbf{Lugar} & \textbf{Sintaxis} \\
\midrule
Espacial, radio & predicado del \texttt{WHERE} & \texttt{col WITHIN\_RADIUS((lon, lat), r)} \\
Espacial, k-NN & cláusula \texttt{ORDER BY} & \texttt{ORDER BY DISTANCE(col, (lon, lat))} \\
Textual & predicado del \texttt{WHERE} & \texttt{col MATCH 'términos'} \\
Textual, puntaje & lista de proyección & \texttt{COSINE\_SCORE() AS rel} \\
Vectorial & lista de proyección & \texttt{col \textless=\textgreater{} '[v1, ..., vn]' AS sim} \\
\bottomrule
\end{tabularx}
\end{table}

El radio es un predicado porque recorta el conjunto de filas y se puede combinar
con \texttt{AND} como cualquier otra condición. El k-NN, en cambio, no recorta
nada: ordena, y el recorte lo hace el \texttt{LIMIT}. Por eso vive en el
\texttt{ORDER BY} y no en el \texttt{WHERE}, que fue la primera cosa que tuvimos
que decidir y la que más ordenó el resto del diseño.

La distancia vectorial es el caso más incómodo de los cinco. No es un predicado
ni una cláusula de orden: es una columna calculada que no está en la tabla. El
operador \texttt{\textless=\textgreater{}} aparece en la proyección, produce una similitud por fila y
obliga a declarar un alias, porque sin nombre el \texttt{ORDER BY} no tiene cómo
referirse a ella. El motor rechaza la consulta si falta:

\begin{lstlisting}[language=SQL, caption={Declaración y uso del índice vectorial}]
CREATE TABLE rostros (
  id        INT PRIMARY KEY,
  lote      CHAR(8),
  embedding VARCHAR(23)
) USING HEAP;

CREATE INDEX ix_rostros_ivf ON rostros (embedding)
  USING IVF WITH (mode = 'face', nprobe = 4);

SELECT id, lote, embedding <=> '[0.111721, 0.152196, ...]' AS similitud
FROM rostros
WHERE id < 1000
ORDER BY similitud DESC
LIMIT 5;
\end{lstlisting}

\subsubsection{La columna indexada guarda una referencia, no el vector}

La tabla no almacena los 128 flotantes de cada rostro. Guarda la ruta de la
imagen, y es el índice el que resuelve de ahí el vector. Esto no fue una
concesión: es la misma idea con la que el R-Tree del bloque espacial indexa un
\texttt{POINT(lon lat)} escrito en una columna de texto. La tabla guarda la
referencia y la estructura guarda la geometría.

La consecuencia práctica es que un registro de \texttt{rostros} ocupa 36 bytes y
entran 101 por página, así que las 13\,000 filas caben en 129 páginas. Si el
vector viviera en la tabla, cada fila pesaría 512 bytes solo de flotantes y la
misma tabla ocuparía más de 1\,700 páginas, con el agravante de que ninguna
consulta relacional necesita leer esos bytes.

\subsection{Acoplamiento de las Tres Familias de Índice}

\subsubsection{Un contrato que no sabe de geometría ni de texto}

Las tres estructuras nuevas se enchufan por la misma interfaz que ya usaban el
árbol B$^+$ y el hash extensible. Un índice declara qué puede servir y el motor
le pregunta antes de registrarlo; si no puede, devuelve el motivo en texto y ese
motivo es lo que llega al usuario:

\begin{lstlisting}[caption={Motivo de rechazo de un índice}]
CREATE INDEX ix ON rostros (id) USING IVF;

  rostros.id debe guardar la ruta de la imagen en CHAR/VARCHAR
\end{lstlisting}

El despachador mantiene la lista de candidatos y encamina cada operación al que
aceptó ese índice. Lo que no acepta nadie cae a un sustituto en memoria, de modo
que una consulta sigue respondiendo aunque la estructura que le correspondía no
esté disponible. El cliente muestra en qué estructura quedó cada índice, y en
las mediciones de este informe los tres corren sobre la real:

\begin{table}[H]
\centering
\caption{Las tres familias conviviendo en un solo catálogo}
\begin{tabularx}{\textwidth}{@{}llX@{}}
\toprule
\textbf{Familia} & \textbf{Dónde corre} & \textbf{Estructura que la atiende} \\
\midrule
\texttt{RTREE} & en el proceso del motor & R-Tree paginado en disco \\
\texttt{INVERTED} & en el proceso del motor & índice invertido en disco (SPIMI) \\
\texttt{IVF} & en un servicio aparte, por HTTP & IVF del bloque vectorial \\
\bottomrule
\end{tabularx}
\end{table}

\subsubsection{El índice que vive en otro proceso}

El IVF es el único que no corre dentro del motor. Sus vectores los produce un
modelo que arrastra opencv y onnxruntime, y meter eso en la imagen del motor la
multiplicaba de tamaño. El adaptador implementa el mismo contrato que los demás
y traduce por HTTP, con una complicación propia: el servicio identifica sus
rostros con un número correlativo y el motor trabaja con RIDs. El adaptador
construye el mapa entre los dos al crear el índice, leyendo la columna de rutas,
y es lo único que necesita para traducir en ambas direcciones.

Ese mapa es también lo único que se pierde al reiniciar, y encontrarlo nos costó.
Con el catálogo persistido, el motor arrancaba con el índice IVF declarado pero
vacío, así que la consulta vectorial se planificaba bien y devolvía cero filas.
Los índices espacial y textual tenían su rutina de reapertura y el vectorial no.
Ahora el arranque lo reconstruye recorriendo la tabla, que es donde las rutas ya
están guardadas, y hay una prueba que levanta el motor dos veces sobre el mismo
catálogo para fijar ese comportamiento.

\subsection{Planificador Híbrido}

\subsubsection{Las dos rutas}

Cuando una consulta trae un filtro relacional y además un ranking vectorial, hay
dos formas de resolverla y dan el mismo resultado por caminos de costo muy
distinto.

\begin{table}[H]
\centering
\caption{Las dos estrategias de una consulta híbrida}
\begin{tabularx}{\textwidth}{@{}llX@{}}
\toprule
\textbf{Estrategia} & \textbf{Nodo} & \textbf{Qué hace} \\
\midrule
Pre-filtering & \texttt{IVFRefine} & la vía relacional acota las filas y el índice
solo calcula la similitud de esas \\
Post-filtering & \texttt{IVFScan} con un \texttt{Filter} encima & el índice sondea
sus listas y devuelve candidatos ordenados; el filtro se aplica después \\
Ranking directo & \texttt{IVFScan} & no hay filtro relacional \\
\bottomrule
\end{tabularx}
\end{table}

La estrategia que corrió queda escrita en el nodo del plan y se ve tanto en la
consola SQL como en la pantalla de identificación, que es lo que el enunciado
pide reportar.

\subsubsection{La regla de decisión}

La primera versión comparaba cuántas filas deja el filtro contra cuántos
candidatos pediría el índice, y elegía pre-filtering cuando el filtro dejaba
menos. Suena razonable y está mal por un motivo que solo se ve midiendo: el
post-filtering aplica el filtro sobre lo que el índice devolvió, no sobre lo que
el índice miró. Si el filtro es muy selectivo, de los resultados del índice no
sobrevive casi ninguno y la consulta responde con menos filas de las pedidas sin
avisar de nada. Con \texttt{LIMIT 5} y un filtro que deja el 4 por ciento de la
tabla, nos devolvía una o dos.

La regla que quedó mira eso directamente. Si $C$ es el número de filas que el
índice puede entregar, $f$ la fracción de la tabla que pasa el filtro y $k$ las
filas pedidas, el post-filtering solo sirve si
\begin{equation}
    \left\lfloor C \cdot f \right\rfloor \ge m \cdot k ,
\end{equation}
con $m = 2$ de margen, porque cuántas sobreviven es un sorteo alrededor de
$C \cdot f$ y pedir lo justo falla la mitad de las veces. Cuando la condición no
se cumple, el motor acota antes aunque eso cueste recorrer la tabla entera.
Elegido el post-filtering, el nodo no pide $k$ filas al índice sino
\begin{equation}
    \min\left(C,\; m \cdot \left\lceil k / f \right\rceil\right),
\end{equation}
de manera que después del filtro sigan quedando las $k$. El plan deja ese número
a la vista en el campo \texttt{pide}.

Hay un tercer caso que al principio no habíamos previsto. Si el planificador no
puede estimar la selectividad del filtro, porque es una igualdad sobre una
columna sin estadísticas, no sabe si el post-filtering alcanza. Dar por bueno que
el filtro no reduce nada lleva a elegirlo y a devolver menos filas; por eso ante
una estimación que no aporta información el motor prefiere pre-filtering, que es
correcto siempre. La diferencia se nota en el filtro \texttt{lote = '00000'} de
nuestras pruebas, que deja mil filas y el catálogo no sabe cuántas.

\begin{lstlisting}[caption={Los dos planes para la misma foto de consulta}]
WHERE id < 100                          ->  pre-filtering
-> Project (columns=id, lote, similitud)
  -> Limit (rows=5)
    -> IVFRefine (index=ix_rostros_ivf dim=128 top_k=5)
      -> Filter (condition=id < 100)
        -> SeqScan (table=rostros)

WHERE id > 2                            ->  post-filtering
-> Project (columns=id, lote, similitud)
  -> Limit (rows=5)
    -> Filter (condition=id > 2)
      -> IVFScan (index=ix_rostros_ivf dim=128 top_k=5 pide=12 nprobe=4)
\end{lstlisting}

\subsubsection{Lo que solo apareció con el dataset completo}

Tres errores de esta parte sobrevivieron a las pruebas unitarias y aparecieron
recién al cargar los 13\,000 rostros. Los tres tenían la misma causa de fondo:
las pruebas corren sobre tablas de cinco filas, y a esa escala cualquier
constante parece razonable.

El primero fue que la decisión leía las filas que reportaba el plan relacional.
Un \texttt{SeqScan} con un \texttt{Filter} encima sigue informando que devuelve
la tabla entera, porque el filtro no cambia su estimación, así que el
planificador creía que ningún filtro reducía nada y elegía post-filtering
siempre. La selectividad hay que calcularla aparte, sobre los rangos sargables
del predicado.

El segundo fue la estimación de cuánto mira el índice. Estaba fijada en cuatro
veces $k$, que son veinte vectores para una consulta de cinco, cuando el IVF con
\texttt{nprobe = 4} sobre 128 listas sondea 1\,340 de 13\,000. Un orden de
magnitud de diferencia, y con el número chico el pre-filtering no ganaba nunca.
Ahora el servicio informa cuántas listas tiene su índice y la estimación sale de
ahí, de la fracción $\text{nprobe}/\text{listas}$, que no depende de $k$.

El tercero fue el tope de resultados. El servicio vectorial no devuelve más de
50 filas por consulta, un límite sensato para una interfaz de búsqueda y un techo
duro para el post-filtering. Pedir 326 candidatos no servía de nada y seguíamos
recibiendo dos filas. El planificador ahora conoce ese tope y lo usa como la $C$
de la fórmula, que es lo que movió la frontera a donde las mediciones la
encuentran.

\subsection{Aplicación de Identificación Facial}

La aplicación recibe una foto y responde con los rostros más parecidos del
conjunto indexado, pero no va al índice por su cuenta: construye una sentencia
SQL con el vector como literal y la ejecuta por el mismo camino que cualquier
otra consulta. Esa decisión tiene una consecuencia que vale el esfuerzo: la
pantalla puede mostrar el plan que corrió, los bloques que leyó y la estrategia
que el planificador eligió, porque son los mismos datos que devuelve la consola
SQL.

\begin{table}[H]
\centering
\caption{Endpoints añadidos al motor}
\begin{tabularx}{\textwidth}{@{}llX@{}}
\toprule
\textbf{Método} & \textbf{Ruta} & \textbf{Función} \\
\midrule
\texttt{GET} & \texttt{/api/identificacion/estado} & Si el bloque vectorial
responde y cuántos rostros tiene indexados \\
\texttt{POST} & \texttt{/api/identificacion/buscar} & Recibe la foto, arma la
consulta y devuelve filas, plan, estrategia y métricas \\
\texttt{GET} & \texttt{/api/identificacion/imagen} & Sirve una imagen indexada
por su ruta \\
\bottomrule
\end{tabularx}
\end{table}

El último endpoint existe por una razón de navegador. La galería de resultados
necesita mostrar las imágenes, que viven en el servicio vectorial en otro
puerto, y pedirlas desde ahí las bloqueaba por origen cruzado. El motor las
reenvía, con lo que el cliente habla con un solo origen igual que en el primer
entregable.

La pantalla tiene tres paneles: la foto de consulta con el filtro relacional y el
número de resultados, la galería ordenada por similitud con el primer puesto
marcado, y el panel que explica cómo se resolvió, con la sentencia que se
ejecutó, el reparto de tiempos y la ruta de acceso. Cambiar el filtro y ver cómo
el sello de la estrategia pasa de una a otra es la demostración más directa de
que el planificador está decidiendo de verdad.

\begin{table}[H]
\centering
\caption{Reparto de una identificación completa, sin filtro relacional}
\begin{tabular}{@{}lrl@{}}
\toprule
\textbf{Etapa} & \textbf{ms} & \textbf{Dónde ocurre} \\
\midrule
Vector de la foto (YuNet + SFace) & 48{,}50 & servicio vectorial \\
Análisis léxico y sintáctico & 0{,}39 & motor \\
Planificación & 0{,}06 & motor \\
Ejecución & 4{,}97 & motor y servicio \\
\midrule
\textbf{Total} & \textbf{53{,}92} & \\
\bottomrule
\end{tabular}
\end{table}

Nueve de cada diez milisegundos se los lleva calcular el vector de la foto, que
es detección más una pasada de la red de reconocimiento sobre una imagen de
1024 por 1024. Todo lo que hace el motor, incluido parsear una sentencia que
lleva 128 flotantes escritos como literal, suma cinco milisegundos. La
planificación cuesta 0{,}06, de modo que el trabajo que describe este informe no
es el que domina el tiempo de respuesta, y eso está bien: significa que la
consulta híbrida sale prácticamente gratis frente al costo de mirar la foto.

Sobre la calidad del resultado, la foto de consulta se recupera a sí misma con
similitud 1{,}000 y el segundo puesto queda en 0{,}373. Esa separación es la que
hace legible la galería: el primer rostro es el mismo y los siguientes son
parecidos de lejos.

\subsection{Experimento 9: Estrategia de Filtrado en Consultas Híbridas}

\begin{lstlisting}[language=bash, caption={Ejecución del experimento}]
python benchmarks/experimento9_hibrido.py
\end{lstlisting}

El arnés monta la tabla de 13\,000 rostros en un directorio temporal con el
backend de disco, toma una foto del propio conjunto como consulta y repite cada
medición cinco veces después de una corrida de calentamiento, quedándose con la
mediana. Deja tres archivos CSV en \texttt{benchmarks/results/}. Las cifras de
bloques son las que el \texttt{DiskCounter} registró.

\begin{table}[H]
\centering
\caption{Estrategia elegida según la selectividad del filtro, con $k = 5$}
\begin{tabular}{@{}lrrlrrr@{}}
\toprule
\textbf{Filtro} & \textbf{Filas} & \textbf{\% tabla} & \textbf{Estrategia} &
\textbf{Pide} & \textbf{Bloques} & \textbf{ms} \\
\midrule
sin filtro & 13\,000 & 100{,}00 & directo & 5 & 5 & 1{,}63 \\
\texttt{id \textless{} 10} & 10 & 0{,}08 & pre & 5 & 128 & 23{,}47 \\
\texttt{id \textless{} 100} & 100 & 0{,}77 & pre & 5 & 128 & 23{,}82 \\
\texttt{id \textless{} 500} & 500 & 3{,}85 & pre & 5 & 128 & 24{,}60 \\
\texttt{id \textless{} 1000} & 1\,000 & 7{,}69 & pre & 5 & 128 & 26{,}25 \\
\texttt{id \textless{} 2600} & 2\,600 & 20{,}00 & pre & 5 & 128 & 30{,}46 \\
\texttt{id \textless{} 5000} & 5\,000 & 38{,}46 & post & 28 & 28 & 1{,}79 \\
\texttt{id \textless{} 9000} & 9\,000 & 69{,}23 & post & 16 & 16 & 1{,}24 \\
\texttt{id \textgreater{} 2} & 12\,997 & 99{,}98 & post & 12 & 12 & 1{,}11 \\
\bottomrule
\end{tabular}
\end{table}

Las nueve consultas devolvieron las cinco filas pedidas, que es la invariante
que la regla de decisión y la sobre-emisión existen para sostener.

\begin{figure}[H]
\centering
\begin{tikzpicture}
\begin{axis}[figura,
    height=6.8cm,
    xmode=log, ymode=log,
    xlabel={Fracción de la tabla que deja pasar el filtro (\%)},
    ylabel={Bloques leídos},
    xmin=0.05, xmax=200, ymin=6, ymax=330,
    xtick={0.1,1,10,100},
    xticklabels={0{,}1\%,1\%,10\%,100\%},
    ytick={10,30,100,128},
    yticklabels={10,30,100,128},
    legend pos=south west,
]
% La referencia y la marca del cambio van antes que las series, para quedar
% por debajo de ellas.
\addplot[cHeap, dashed, mark=none, forget plot] coordinates {(0.05,128)(200,128)};
\node[nota, anchor=south west, text=cHeap] at (axis cs:0.055,133)
  {escaneo completo de la tabla};

\draw[cNota, dotted, line width=0.6pt] (axis cs:28,6) -- (axis cs:28,105);
\node[nota, anchor=east] at (axis cs:24,46)
  {el índice ya puede entregar\\las 5 filas tras el filtro};

\addplot[cBtree, mark=*]
  coordinates {(0.08,128)(0.77,128)(3.85,128)(7.69,128)(20,128)};
\addlegendentry{Pre-filtering (IVFRefine)}
\addplot[cSeq, mark=triangle*]
  coordinates {(38.46,28)(69.23,16)(99.98,12)};
\addlegendentry{Post-filtering (IVFScan)}
\end{axis}
\end{tikzpicture}
\caption{Bloques leídos según lo selectivo que sea el filtro relacional, con
$k = 5$ y ambos ejes logarítmicos. El pre-filtering paga el escaneo completo de
la tabla sin importar cuántas filas sobrevivan, porque \texttt{rostros} no tiene
índice sobre la columna del filtro; el post-filtering solo lee las filas que
termina devolviendo. La línea punteada vertical marca dónde el planificador
cambia de estrategia, que es donde el índice empieza a poder entregar las cinco
filas pedidas una vez aplicado el filtro.}
\label{fig:e9-estrategia}
\end{figure}

El salto entre las dos estrategias es de dos órdenes de magnitud en tiempo, de
30 milisegundos a 1{,}8, y de más de cuatro veces en bloques. La razón está en
que \texttt{rostros} no tiene índice sobre \texttt{id}: el pre-filtering acota
las filas recorriendo la tabla completa, así que su costo es el del escaneo
secuencial más el cálculo de similitud de las filas que pasaron. Es el precio de
poder responder la consulta completa, y el planificador lo paga solo cuando la
otra vía no puede.

\begin{table}[H]
\centering
\caption{Los dos extremos entre los que se mueven las estrategias}
\begin{tabular}{@{}lrr@{}}
\toprule
\textbf{Ruta} & \textbf{Bloques} & \textbf{ms} \\
\midrule
Escaneo completo con filtro (\texttt{SeqScan} + \texttt{Filter}) & 128 & 21{,}37 \\
Ranking por el IVF, sin filtro & 5 & 1{,}63 \\
\bottomrule
\end{tabular}
\end{table}

La tercera parte del experimento varía \texttt{nprobe}, que es cuántas de las
128 listas del índice se sondean. Lo interesante es que mueve mucho el trabajo
del índice y nada la estrategia:

\begin{table}[H]
\centering
\caption{Efecto de \texttt{nprobe} sobre el sondeo y sobre la decisión}
\begin{tabular}{@{}rrrrlr@{}}
\toprule
\textbf{nprobe} & \textbf{Candidatos} & \textbf{\% colección} &
\textbf{ms sin filtro} & \textbf{Estrategia con \texttt{id \textless{} 1000}} &
\textbf{Bloques} \\
\midrule
1 & 353 & 2{,}72 & 1{,}13 & pre & 128 \\
2 & 607 & 4{,}67 & 1{,}17 & pre & 128 \\
4 & 1\,340 & 10{,}31 & 1{,}22 & pre & 128 \\
8 & 2\,345 & 18{,}04 & 1{,}23 & pre & 128 \\
16 & 4\,474 & 34{,}42 & 1{,}59 & pre & 128 \\
\bottomrule
\end{tabular}
\end{table}

Los candidatos crecen de 353 a 4\,474, trece veces, y el tiempo sube de 1{,}13 a
1{,}59 milisegundos. La comparación de vectores está vectorizada en el servicio y
cuesta poco frente al resto, así que subir \texttt{nprobe} sale casi gratis en
tiempo aunque multiplique el trabajo nominal. La estrategia, en cambio, no se
mueve: con \texttt{id \textless{} 1000} el post-filtering seguiría sin alcanzar porque el
techo no es cuánto sondea el índice sino cuántas filas accede a devolver, que son
50. Sondear más no ayuda si el resultado se recorta igual.

Esa es la lección del experimento, y no la esperábamos. Pensábamos que la
decisión entre pre y post filtering iba a depender del costo de cada ruta, y
depende de algo anterior: de si la ruta puede responder la pregunta. El costo
solo decide entre las rutas que sí pueden, y en esta tabla, sin un índice sobre
la columna del filtro, cuando el post-filtering puede es además el más barato por
un margen amplio.

\subsection{Verificación}

La parte del motor tiene 248 pruebas automatizadas, de las cuales 12 cubren el
camino vectorial con un doble del servicio, de modo que corren sin opencv ni
modelos. Además hay una comprobación de extremo a extremo que levanta una tabla
por bloque, corre una consulta de cada familia y revisa que el plan elegido sea
el que corresponde:

\begin{lstlisting}[language=bash, caption={Comprobación de los cuatro bloques}]
python scripts/verificar_p2.py
\end{lstlisting}

\begin{lstlisting}[caption={Salida de la comprobación}]
Bloque A  espacial (R-Tree paginado)
  ok   busqueda por radio usa el R-Tree               3 locales
  ok   el radio deja fuera lo que no alcanza          ids [1, 2, 3]
  ok   k-NN usa el R-Tree                             2 < 3 < 1
Bloque B  textual (indice invertido SPIMI)
  ok   ranking por coseno usa el indice invertido     2 documentos
  ok   consulta booleana con NOT                      ids [5]
Bloque C  vectorial (IVF sobre rostros) y planificador hibrido
  ok   el motor ve el servicio vectorial              13000 imagenes, 12998 rostros
  ok   la foto de consulta se encuentra a si misma    00000/00000.png -> 100.0%
  ok   un filtro selectivo elige pre-filtering
  ok   un filtro debil elige post-filtering
Bloque D  catalogo y vertebra del motor
  ok   las tres familias de indice conviven en un catalogo  INVERTED IVF RTREE
  ok   los indices corren sobre su estructura real, no el sustituto
  ok   la via vectorial lee menos bloques que el escaneo completo  5 contra 129
\end{lstlisting}

El conjunto indexado son las carpetas \texttt{00000} a \texttt{13000} de FFHQ,
13\,000 imágenes de 1024 por 1024 que pesan 17 GB. El detector encontró rostro en
12\,998, y las dos que faltan son fotos donde la cara queda cortada o muy de
perfil.
